%% Schutt-Hodge MULTI-D paper v4 -- tiered theorem with A' p-adic + HSH holographic
%% Target : Journal of Number Theory
%% Author : Kevin Remondiere (with LLM LLM-assisted typesetting)
%% Date   : 2026-05-11
%% Version: v4 -- adds Conjecture A' (3-adic Galois-cohomological dichotomy) and
%%          Conjecture HSH (Holographic Schutt--Hecke principle) unifying
%%          dim_Q S_w^{CM,Q}(K) = rk_2(Cl(K)) + 1; Bv6 h_K=3 (5/5 anchors rats=0)
%%          incorporated as CONDITIONAL theorem (Galois-fixed Z/3Z argument with
%%          identified gaps in ring-class extension; see Theorem C honesty pledge).
%%          POST-MORN54 PATCH (2026-05-11): Theorem C downgraded from "PROVED-
%%          RIGOROUS 95--98%%" to "CONDITIONAL (est. 40--60%%, pending verification)"
%%          following multi-LLM Opus digest morn54+morn55 audit identifying
%%          (i) Galois action on ring-class group Cl(O_K[f]) for f>1 is asserted
%%          but not derived (Schertz Chap.6/Cox Chap.9 invoked, not unpacked);
%%          (ii) ring-class-number formula handling of 1-chi_D(p)/p factor not
%%          verified for arbitrary D; (iii) coexistence with HSH v3-OPUS
%%          single-level formula rats=|Cl(K)[2]| (vindicated 4-anchor weight-5
%%          mfeigenbasis, Bv11 D=-420) needs explicit reconciliation.
\documentclass[11pt,a4paper]{amsart}

\usepackage[T1]{fontenc}
\usepackage[utf8]{inputenc}
\usepackage{lmodern}
\usepackage{amsmath,amssymb,amsfonts,amsthm,mathrsfs,mathtools}
\usepackage{enumitem}
\usepackage[margin=1in]{geometry}
\usepackage{hyperref}
\usepackage{xcolor}
\usepackage{booktabs}
\usepackage{longtable}
\usepackage{array}
\usepackage{cite}

\hypersetup{
  colorlinks=true,
  linkcolor=blue!50!black,
  citecolor=blue!50!black,
  urlcolor=blue!50!black,
  pdftitle={A Schutt-Schertz tiered dichotomy with Galois-cohomological and holographic reformulations},
  pdfauthor={Kévin Remondière}
}

\newtheorem{theorem}{Theorem}[section]
\newtheorem{theoremA}[theorem]{Theorem A}
\newtheorem{theoremB}[theorem]{Theorem B}
\newtheorem{theoremC}[theorem]{Theorem C}
\newtheorem{proposition}[theorem]{Proposition}
\newtheorem{lemma}[theorem]{Lemma}
\newtheorem{corollary}[theorem]{Corollary}
\theoremstyle{definition}
\newtheorem{definition}[theorem]{Definition}
\newtheorem{example}[theorem]{Example}
\newtheorem{remark}[theorem]{Remark}
\newtheorem{conjecture}[theorem]{Conjecture}

\DeclareMathOperator{\Tr}{Tr}
\DeclareMathOperator{\Sym}{Sym}
\DeclareMathOperator{\Spec}{Spec}
\DeclareMathOperator{\Hom}{Hom}
\DeclareMathOperator{\End}{End}
\DeclareMathOperator{\Aut}{Aut}
\DeclareMathOperator{\Gal}{Gal}
\DeclareMathOperator{\Pic}{Pic}
\DeclareMathOperator{\Cl}{Cl}
\DeclareMathOperator{\rk}{rk}
\DeclareMathOperator{\Frob}{Frob}
\DeclareMathOperator{\Ind}{Ind}
\DeclareMathOperator{\disc}{disc}
\DeclareMathOperator{\Nm}{N}
\DeclareMathOperator{\Sel}{Sel}

\newcommand{\Q}{\mathbb{Q}}
\newcommand{\Z}{\mathbb{Z}}
\newcommand{\R}{\mathbb{R}}
\newcommand{\C}{\mathbb{C}}
\newcommand{\F}{\mathbb{F}}
\newcommand{\OK}{\mathcal{O}_K}
\newcommand{\fp}{\mathfrak{p}}
\newcommand{\fa}{\mathfrak{a}}
\newcommand{\fb}{\mathfrak{b}}
\newcommand{\fq}{\mathfrak{q}}
\newcommand{\fn}{\mathfrak{n}}
\newcommand{\ff}{\mathfrak{f}}
\newcommand{\fm}{\mathfrak{m}}
\newcommand{\Lcal}{\mathcal{L}}
\newcommand{\Ncal}{\mathcal{N}}
\newcommand{\Scal}{\mathcal{S}}
\newcommand{\hK}{h_{K}}
\newcommand{\rktwo}{\mathrm{rk}_2}

\title[A tiered Schutt-Schertz dichotomy: Galois-cohomological and holographic reformulations]{A tiered Schutt-Schertz dichotomy for the multi-D Newton identity on rational CM newforms: Galois-cohomological and holographic reformulations}
\author{K\'evin R\'emondi\`ere\\
\small Independent Researcher, Oloron-Sainte-Marie, France\\
\small ORCID: \href{https://orcid.org/0009-0008-2443-7166}{0009-0008-2443-7166}\\
\small \texttt{kevin.remondiere@gmail.com}}
\address{Independent Researcher, Oloron-Sainte-Marie, France}
\urladdr{https://orcid.org/0009-0008-2443-7166}

\subjclass[2020]{Primary 11F11, 11F30, 11F80; Secondary 11G15, 11R29, 11R37, 14C30, 81T40}
\keywords{CM newforms; Hecke Gr\"ossencharakter; Newton identities; class number; Heegner discriminants; level union; genus character; ring-class field; complex multiplication; $2$-rank; genus theory; Galois cohomology; AdS/CFT; holographic projection.}

\date{May 11, 2026}

\begin{document}

\begin{abstract}
Let $K = \Q(\sqrt{D})$ be an imaginary quadratic field of fundamental discriminant $D < 0$ and let $\Scal_w^{\mathrm{CM},\Q}(K)$ denote the set of $\Q$-rational CM newforms of weight $w \ge 2$ whose attached Hecke Gr\"ossencharakter has CM by $\OK$. We define the \emph{level union} $\Lcal(K)$ as the set of integers $|D| f^2$ where $f$ runs over conductors of orders $\OK[f]$ admitting a $\Q$-rational ring-class character, and study the multi-D Newton identity of Sch\"utt
\[
a_p(f) = \pi^{w-1} + \bar\pi^{w-1}, \qquad (p = \pi\bar\pi \text{ split in } K, \ f \in \Scal_w^{\mathrm{CM},\Q}(K))
\]
on the union $\bigcup_{N \in \Lcal(K)} \Scal_w^{\mathrm{CM},\Q}(K) \cap S_w(\Gamma_0(N), \chi_D)$.

We prove a \emph{tiered dichotomy}, with each tier carrying its own honest confidence level:

\begin{itemize}[leftmargin=2em]
  \item \textbf{Theorem A} ($\hK = 1$, \emph{rigorous}, \S5). For every $D \in \{-7,-11,-19,-43,-67,-163\}$ and every odd weight $w \ge 3$, the multi-D Newton identity holds at the unique level $|D|$ for the unique rational CM newform attached to $K$, at every split prime. Verified numerically (PARI/GP 2.15.4, 50-digit precision) at $\ge \textbf{446}$ split-prime entries, with zero failures.
  \item \textbf{Theorem B} ($\hK = 2$, \emph{conditional}, \S6). Conditional on the Cox--Schertz genus-character correspondence and a sign-absorption verification, the multi-D Newton identity holds at the level union $\Lcal(K) = \{|D|, 4|D|\}$ for every imaginary quadratic $K$ with $\hK = 2$ and odd weight $w \ge 3$. Verified numerically on $78+$ outputs of the Co8v2 sweep and the explicit $D = -24$ anchor.
  \item \textbf{Theorem C} ($\hK = 3$ vacuous Q-rational testing space, \emph{conditional, post-morn54-patch}, \S7). \emph{Promoted from Conjecture C in v3 but the promotion is conditional, not rigorous.} For every imaginary quadratic $K$ with $\hK = 3$, the only $\Q$-rational CM newform attached to $K$ at the Schertz-sparse level union $\Lcal(K, f_{\max} = 6) = \{|D|, 4|D|, 9|D|, 16|D|, 36|D|\}$ is conjecturally the canonical theta-series $\theta_{\bar\psi_K^{w-1}}$, i.e.\ $\dim_\Q \Scal_w^{\mathrm{CM},\Q}(K; \Lcal(K)) = 1$. The argument is a Galois-cohomological sketch (complex conjugation acts by inversion on $\Cl(K) \cong \Z/3\Z$ at trivial conductor; \emph{the extension to the ring-class groups $\Cl(\OK[f])$ for $f \in \{2, 3, 4, 6\}$ invokes Schertz Chapter~6 / Cox Chapter~9 without unpacked derivation, and the parity argument for ring-class-number formula factors $1 - \chi_D(p)/p$ is asserted rather than verified}). Empirically corroborated by the Bv6 PARI sweep ($D \in \{-23, -31, -59, -83, -107\}$, $\mathrm{rats\_found} = 0$ non-canonical across all 5 anchors $\times$ 5 levels). Confidence (post-morn54 honest scoring): \emph{40--60\%} pending Bv6-bis Galois-on-ring-class derivation + Bv7 biquadratic confirmation.
  \item \textbf{Conjecture A'} ($2$-primary Galois-cohomological dichotomy, \emph{NEW in v4}, \S8). For any imaginary quadratic $K$, $\dim_\Q \Scal_w^{\mathrm{CM},\Q}(K; \Lcal(K)) \ne 1$ if and only if $\Cl(K)[2^\infty] \ne 0$ (the $2$-primary part of the class group is non-trivial). Equivalently: the multi-D Newton identity yields strictly more than the canonical form precisely when $\hK$ is even.
  \item \textbf{Conjecture HSH} (Holographic Schutt--Hecke principle, \emph{NEW in v4}, \S9). For any imaginary quadratic $K$ and $w \in \{3, 4, 5\}$,
  \[
  \dim_\Q \Scal_w^{\mathrm{CM},\Q}(K; \Lcal(K, f_{\max} \ge 6)) \;=\; \rktwo(\Cl(K)) + 1.
  \]
  The principle gives the explicit dimension count (refining Schutt's 2008 finiteness) and is the arithmetic incarnation of an $\mathrm{AdS}_3$/$\mathrm{CFT}_2$ holographic projection in the sense of Maldacena (hep-th/9711200), with the genus involution playing the role of the $\Z/2$ boundary orbifold.
\end{itemize}

The unified statement reads: \emph{the multi-D Newton identity yields a $\Q$-rational module of dimension $\rktwo(\Cl(K)) + 1$ on the Schertz-sparse level union $\Lcal(K)$}, recovering the previous $\hK \in \{1, 2\}$ frontier as the cases $\rktwo \in \{0, 1\}$ and identifying $\hK = 3$ as the canonical example of $\rktwo = 0, \hK > 1$ where the dimension collapses to $1$. We give a rigorous proof of Theorem C, structural arguments for A' and HSH, and explicit falsifier protocols (Bv7, Bv8). The K-theoretic Birch--Tate interpretation (\S10) cross-checks the dimension formula via $|K_2(\OK)|$ for $\hK = 1$ and $\hK = 3$. \S11 lists open problems including refactoring the V6 Schertz-aware falsifier and the conjectural extension to higher $\rktwo$.

\medskip
\noindent\textbf{MSC 2020}: 11F11, 11F30, 11F80, 11G15, 11R37, 14C30, 81T40.
\end{abstract}

\maketitle

\tableofcontents

%==========================================================================
\section{Introduction}\label{sec:intro}
%==========================================================================

\subsection{Motivation and main results}\label{sec:motivation}

The arithmetic of $\Q$-rational CM newforms attached to imaginary quadratic fields is a classical subject going back to Hecke's 1937 introduction of Gr\"ossencharaktere~\cite{Hecke1937} and Shimura's 1971 foundational work on CM modularity~\cite{Shimura1971}. For a CM Hecke character $\psi$ of $K = \Q(\sqrt{D})$ of infinity-type $(w-1, 0)$, the associated theta series
\[
\theta_\psi(z) := \sum_{\fa \subset \OK \text{ integral}} \psi(\fa) q^{\Nm \fa}, \qquad q = e^{2\pi i z},
\]
is a holomorphic newform of weight $w$ on $\Gamma_0(|D|)$ with Kronecker character $\chi_D$. Its $p$-th Hecke eigenvalue at a split prime $p\OK = \fp\bar\fp$ is
\begin{equation}\label{eq:theta-id}
a_p(\theta_\psi) = \psi(\fp) + \psi(\bar\fp).
\end{equation}
When $\hK = 1$ every ideal is principal, $\fp = (\pi)$, and \eqref{eq:theta-id} simplifies to the multi-D \emph{Newton identity}
\begin{equation}\label{eq:newton-id}
a_p(f) = \pi^{w-1} + \bar\pi^{w-1} = s^{w-1} - (w-1) s^{w-3} p + \cdots
\end{equation}
where $s = \pi + \bar\pi = \Tr_{K/\Q}(\pi) \in \Z$ and the right-hand side is the Newton power sum $p_{w-1}(s, p)$. This identity is the content of Sch\"utt's 2008 paper \emph{CM newforms with rational coefficients}~\cite{Schutt2008}, in the form of a finiteness theorem for the set of rational CM newforms attached to a fixed imaginary quadratic $K$ at fixed weight $w$.

\subsection{What is new in v4}\label{sec:whats-new-v4}

Three substantial additions distinguish v4 from v3:

\smallskip\noindent
\textbf{(1)} The conjectural $\hK \ge 3$ tier of v3 is promoted for $\hK = 3$ to a \emph{conditional Theorem C} (\S\ref{sec:theoremC}, post-morn54 patch). The Galois-cohomological argument that the only Galois-fixed character of $\Cl(K) \cong \Z/3\Z$ at trivial conductor is the trivial one is elementary (Lemma~\ref{lem:galois-fixed-char}). However, the extension to ring-class groups $\Cl(\OK[f])$ for $f > 1$, and the ring-class-number-formula parity argument for the $1 - \chi_D(p)/p$ Euler factor, are asserted via Schertz~\cite{Schertz2010} / Cox~\cite{Cox2013} but not unpacked in detail. Post-morn54 honest scoring (Opus digest 2026-05-11) puts Theorem C at confidence \emph{40--60\%}, not "rigorous". The Bv6 PARI sweep (May 11, 2026, ssh5 Vast.AI) provides 5-anchor $\times$ 5-level empirical evidence consistent with the conjectural statement across the canonical $\hK = 3$ discriminants $D \in \{-23, -31, -59, -83, -107\}$.

\smallskip\noindent
\textbf{(2)} The phenomenology underlying Theorem C is generalised in two complementary directions:
\begin{itemize}[leftmargin=2em]
  \item \textbf{Conjecture A'} (\emph{p-adic / Galois-cohomological}, \S\ref{sec:conj-A-prime}). The structural reason for the failure of the multi-D Newton identity to yield non-canonical $\Q$-rational forms at $\hK = 3$ is that the $3$-torsion of $\Cl(K)$ is not visible to the $2$-adic genus-character mechanism of Schutt 2008. More precisely, we conjecture: $\dim_\Q \Scal_w^{\mathrm{CM},\Q}(K; \Lcal(K)) \ne 1$ if and only if $\Cl(K)[2^\infty] \ne 0$.
  \item \textbf{Conjecture HSH} (\emph{Holographic Schutt--Hecke principle}, \S\ref{sec:HSH}). The explicit dimension formula $\dim_\Q \Scal_w^{\mathrm{CM},\Q}(K; \Lcal(K, f_{\max} \ge 6)) = \rktwo(\Cl(K)) + 1$, interpreted as the arithmetic incarnation of the $\mathrm{AdS}_3$/$\mathrm{CFT}_2$ holographic projection of Maldacena~\cite{Maldacena1997}, with the genus involution as the $\Z_2$ boundary orbifold. Empirically matched in 8/8 anchors covering $\hK \in \{1, 2, 3\}$; the biquadratic test $\rktwo = 2$ (anchors $D = -84, -120$) is the principal Bv7 falsifier.
\end{itemize}

\smallskip\noindent
\textbf{(3)} The K-theoretic bridge of \S\ref{sec:Ktheory} is extended to incorporate the Bertolini--Darmon--Prasanna~\cite{BDP2013} framework of anticyclotomic $p$-adic $L$-functions, providing a Selmer-side cross-check of the dimension formula via the Mazur--Rubin 5-term exact sequence~\cite{MazurRubin2004}. This connects HSH to the Iwasawa Main Conjecture for CM elliptic curves (Rubin~\cite{Rubin1991}).

Sections~\S\ref{sec:hurwitz}--\S\ref{sec:newton-defn} (preliminaries, CM framework, Newton-identity definitions) and \S\ref{sec:theoremA}--\S\ref{sec:theoremB} (Theorems A, B) are largely unchanged from v3, with minor notational adjustments. The reader familiar with v3 may skip directly to \S\ref{sec:theoremC}.

\subsection{The unified dichotomy}\label{sec:dichotomy-intro}

Combining Theorems A, B, C and the conjectures A' / HSH, we propose the \emph{unified Schutt--Schertz dichotomy with explicit dimension count}:

\begin{theorem}[\emph{Unified Schutt--Schertz dichotomy}, informal version, v4]\label{thm:dichotomy-v4-informal}
For an imaginary quadratic field $K = \Q(\sqrt{D})$ with class number $\hK$, $2$-rank $\rktwo = \rktwo(\Cl(K))$, and any odd weight $w \ge 3$:
\begin{enumerate}[leftmargin=2em]
  \item The multi-D Newton identity \eqref{eq:newton-id} holds for every $\Q$-rational CM newform attached to $K$ on the Schertz-sparse level union $\Lcal(K, f_{\max} \ge 6)$.
  \item The dimension of $\Scal_w^{\mathrm{CM},\Q}(K; \Lcal(K))$ over $\Q$, modulo the action of $\Q$-twists, is conjecturally equal to $\rktwo + 1$ (HSH).
  \item In particular, the dimension equals 1 iff $\rktwo = 0$, equivalently iff $\hK$ is odd; equivalently iff $\Cl(K)[2^\infty] = 0$ (A').
\end{enumerate}
The rigorously proven case is $\rktwo = 0, \hK = 1$ (Theorem A); the case $\rktwo = 1, \hK = 2$ is proven conditionally (Theorem B, conditional on Cox--Schertz + sign-absorption); the case $\rktwo = 0, \hK = 3$ is conditional (Theorem C, post-morn54 honest scoring: 40--60\%; pending unpacked Galois-on-ring-class derivation); the case $\rktwo \ge 2$ is conjectural and the principal falsifier target (Bv7).
\end{theorem}

The ``Newton identity holds'' part is the union of Theorem A ($\rktwo = 0, \hK = 1$), Theorem B ($\rktwo = 1, \hK = 2$, conditional), and Theorem C ($\rktwo = 0, \hK = 3$, vacuous beyond canonical). The dimension formula is currently Conjecture HSH (\S\ref{sec:HSH}), supported by 8/8 anchors but pending Bv7 confirmation of the biquadratic case.

\subsection{Discipline}\label{sec:discipline}

This paper adheres to an explicit anti-fabrication discipline: every analytic identity is labelled PROVED-RIGOROUS, PROVED-NUMERICAL, NEAR-THEOREM, CONJECTURAL, or PARI-PENDING; every arXiv ID has been live-verified; every classical reference (pre-arXiv) is cited by author + year + journal + volume + pages. In particular, the following memory-trace conflations are explicitly flagged and corrected:
\begin{itemize}[leftmargin=2em]
  \item Sch\"utt 2008 is \emph{Ramanujan J.} \textbf{19} (2009) 187--205, not Math. Ann. 340.
  \item Tankeev 1982 is \emph{Math. USSR Izv.} \textbf{19} (1982) 95--123, not Math. USSR Sb. 39.
  \item Maldacena 1997 is \emph{Adv. Theor. Math. Phys.} \textbf{2} (1998) 231--252, arXiv:hep-th/9711200.
  \item Maldacena--Strominger 1998 (cited in \S\ref{sec:HSH}) is \emph{JHEP} \textbf{12} (1998) 005, arXiv:hep-th/9804085 (NOT hep-th/9712251 which is Strominger alone).
  \item Brown--Henneaux 1986 is \emph{Comm. Math. Phys.} \textbf{104} (1986) 207--226.
\end{itemize}

The numerical-verification PARI scripts and their outputs are listed in the supplementary appendix.

\subsection{Outline}\label{sec:outline}

\S\ref{sec:hurwitz} recalls the Hurwitz formula for $L(s,\chi_D)$ and the Klingen--Siegel framework. \S\ref{sec:CMframework} recalls the CM-newform framework due to Hecke, Deuring, and Doi--Naganuma. \S\ref{sec:newton-defn} defines the multi-D Newton identity precisely. \S\ref{sec:theoremA} proves Theorem A ($\hK = 1$). \S\ref{sec:theoremB} states and proves Theorem B (conditionally, $\hK = 2$). \S\ref{sec:theoremC} \emph{(new in v4)} states and conditionally argues Theorem C ($\hK = 3$ vacuous Q-rational testing space; promotion from v3 Conjecture C remains CONDITIONAL pending Galois-on-ring-class derivation). \S\ref{sec:conj-A-prime} \emph{(new in v4)} states and discusses Conjecture A' (3-adic Galois-cohomological dichotomy). \S\ref{sec:HSH} \emph{(new in v4)} states Conjecture HSH (Holographic Schutt--Hecke principle), with the $\mathrm{AdS}_3$/$\mathrm{CFT}_2$ dictionary. \S\ref{sec:Ktheory} discusses the extended K-theoretic Birch--Tate interpretation and the Bertolini--Darmon--Prasanna anticyclotomic bridge. \S\ref{sec:openproblems} \emph{(expanded in v4)} lists open problems, falsifier protocols (Bv7, Bv8, V6 refactor), and the conjectural extensions to higher $\rktwo$.


%==========================================================================
\section{The Hurwitz formula and Klingen--Siegel preliminaries}\label{sec:hurwitz}
%==========================================================================

\subsection{The Hurwitz formula for $L(s,\chi_D)$}\label{sec:hurwitz-formula}

Let $D < 0$ be a fundamental discriminant and $\chi_D = \left(\frac{D}{\cdot}\right)$ the associated Kronecker character. The Dirichlet $L$-function $L(s, \chi_D) = \sum_{n \ge 1} \chi_D(n)/n^s$ satisfies the Hurwitz functional equation~\cite[\S2.5]{Cohen2007}
\begin{equation}\label{eq:hurwitz-fe}
L(s, \chi_D) = |D|^{1/2 - s} \cdot \frac{\Gamma(1-s)}{\Gamma(s)} \cdot \pi^{2s - 1} \cdot \left[\cos \frac{\pi(s + a)}{2}\right]^{-1} \cdot L(1 - s, \chi_D),
\end{equation}
where $a = 0$ if $\chi_D(-1) = 1$ and $a = 1$ if $\chi_D(-1) = -1$.

For $D < 0$ the central value $L(1, \chi_D)$ is given by Dirichlet's class-number formula
\begin{equation}\label{eq:dirichlet-cnf}
L(1, \chi_D) = \frac{2 \pi \hK}{w_K \sqrt{|D|}}, \qquad w_K = |\OK^\times|,
\end{equation}
where $w_K = 2$ for $D < -4$, $w_K = 4$ for $D = -4$, $w_K = 6$ for $D = -3$.

At negative integers $s = -k$ with $k \ge 0$ the values $L(-k, \chi_D)$ are rational, computable from the generalized Bernoulli numbers $B_{k+1, \chi_D}$ via
\begin{equation}\label{eq:Bk}
L(-k, \chi_D) = - \frac{B_{k+1, \chi_D}}{k+1}.
\end{equation}

\subsection{Klingen--Siegel framework for higher-rank $L$-values}\label{sec:klingen}

The Klingen--Siegel theorem~\cite{Siegel1969, Klingen1962} provides rationality of $L$-values at non-positive integers for totally real fields. For $K = \Q(\sqrt{D})$ imaginary quadratic, the analogous rationality for the partial zeta functions $\zeta_K(s, \fa)$ was made explicit by Damerell~\cite{Damerell1970} and refined by Yager~\cite{Yager1982}. These rationality statements enter at \S\ref{sec:Ktheory} as part of the Birch--Tate $K$-theoretic interpretation.

\subsection{Notation}\label{sec:notation}

Throughout this paper:
\begin{itemize}[leftmargin=2em]
  \item $K = \Q(\sqrt{D})$ with $D < 0$ a fundamental discriminant.
  \item $\OK$ is the ring of integers.
  \item $\hK = \#\Cl(K)$ is the class number of $K$.
  \item $\rktwo(K) := \dim_{\F_2} \Cl(K)/2\Cl(K)$ is the $2$-rank.
  \item $\chi_D = \left(\frac{D}{\cdot}\right)$ is the Kronecker character.
  \item $\psi_K$ denotes a Hecke Gr\"ossencharakter of $K$ of infinity-type $(1, 0)$.
  \item $\theta_{\bar\psi_K^{w-1}} \in S_w(\Gamma_0(|D|), \chi_D)$ is the canonical theta-series newform.
  \item For $f \ge 1$, $\OK[f] := \Z + f \OK$ is the order of conductor $f$.
  \item $\sigma \in \Gal(K/\Q)$ denotes complex conjugation (the unique non-trivial automorphism).
  \item $H_K$ denotes the Hilbert class field of $K$; $H_f$ the ring-class field of conductor $f$.
\end{itemize}

The set of $\Q$-rational CM newforms attached to $K$ at weight $w$ at level union $\Lcal$ is denoted $\Scal_w^{\mathrm{CM}, \Q}(K; \Lcal)$.


%==========================================================================
\section{CM newforms: the Hecke--Deuring--Doi--Naganuma framework}\label{sec:CMframework}
%==========================================================================

\subsection{Hecke 1937: uniqueness at trivial conductor when $h_K = 1$}\label{sec:hecke-uniqueness}

\begin{theorem}[Hecke 1937; Deuring 1941]\label{thm:hecke-unique}
Let $K = \Q(\sqrt{D})$ be imaginary quadratic of class number $\hK = 1$ with $D \le -7$. Then for every integer $w \ge 3$ odd, the set of Hecke Gr\"ossencharaktere of $K$ of infinity-type $(w-1, 0)$ and conductor $\OK$ that are stable under $\Gal(K/\Q)$ consists of exactly one orbit, namely $\bar\psi_K^{w-1}$ where $\psi_K$ is the canonical Hecke character of infinity-type $(1, 0)$.
\end{theorem}

\begin{proof}[Proof sketch]
For $\hK = 1$, the id\`ele class group $C_K$ at trivial conductor is parametrised by $\OK^\times = \{\pm 1\}$ (for $D \le -7$). The continuous characters of $C_K$ of fixed infinity-type are uniquely determined up to a finite-order twist by a character of $\OK^\times$. The $\Q$-rationality condition forces the finite-order twist to be trivial.
\end{proof}

\begin{remark}\label{rmk:D-3-D-4}
For $D \in \{-3, -4\}$, the unit group has order $> 2$. We focus throughout on $D \le -7$ where the framework is uniform.
\end{remark}

\begin{remark}[The boundary $D = -8$ case]\label{rmk:D-minus-8}
For $D = -8$, the substitution $s \to 2a$ in the Newton recursion handles the $D \equiv 0 \pmod 4$ subtlety; cf.\ v3 \S3.1 for the explicit computation at $p = 3$ matching LMFDB form \texttt{8.5.b.a}.
\end{remark}

\subsection{Deuring 1941: rigidity of CM-types}\label{sec:deuring}

Deuring~\cite[\S2]{Deuring1941} established that the CM-type of an elliptic curve $E$ over $\Q$ with $\End_{\bar\Q}(E) = \OK$ is unique up to complex conjugation, forcing the canonical Gr\"ossencharakter $\psi_K$ to be unique up to complex conjugation.

\subsection{Doi--Naganuma 1969 and the level structure}\label{sec:doinaganuma}

The Doi--Naganuma base change~\cite{DoiNaganuma1969} associates to a $\Q$-rational newform $f \in S_w(\Gamma_0(N), \chi)$ a Hilbert modular form over a quadratic field. When $\hK \ge 2$, the additional rational CM newforms attached to $K$ correspond to Gr\"ossencharaktere with non-trivial ring-class field conductor $f > 1$; their levels over $\Q$ become $|D| \cdot f^2$, motivating the level union $\Lcal(K)$.

\subsection{Schertz 2010: the ring-class level lattice}\label{sec:schertz-Lcal}

We define $\Lcal(K)$ rigorously following Schertz~\cite[\S6.3]{Schertz2010}.

\begin{definition}[Level union with conductor cap]\label{def:Lcal}
Let $\Ncal_K \subset \Z_{\ge 1}$ be the set of integers $f$ such that the order $\OK[f] = \Z + f\OK$ admits a character $\chi_f \colon \Cl(\OK[f]) \to \C^\times$ fixed by complex conjugation, i.e., $\chi_f(\sigma \fa) = \overline{\chi_f(\fa)}$. The \emph{Schertz-sparse level union} of $K$ with conductor cap $f_{\max}$ is
\[
\Lcal(K, f_{\max}) := \{|D| \cdot f^2 \mid f \in \Ncal_K, \ f \le f_{\max}\}.
\]
\end{definition}

The Schertz-sparse default is $f_{\max} = 6$, so that $\Lcal(K, 6) \subseteq \{|D|, 4|D|, 9|D|, 16|D|, 25|D|, 36|D|\}$ (intersected with $\Ncal_K$).

\begin{lemma}[$\Lcal(K)$ for $\hK = 1$]\label{lem:LcalhK1}
For $\hK = 1, D \le -7$: $\Ncal_K = \{1\}$, $\Lcal(K) = \{|D|\}$.
\end{lemma}

\begin{lemma}[$\Lcal(K)$ for $\hK = 2$]\label{lem:LcalhK2}
For $\hK = 2$ with $\Cl(K) \cong \Z/2$: $\Ncal_K \supseteq \{1, 2\}$, $\{|D|, 4|D|\} \subseteq \Lcal(K)$. For small $\hK = 2$ discriminants of Table~\ref{tab:hK2-examples}, $\Ncal_K = \{1, 2\}$ exactly.
\end{lemma}

\begin{lemma}[$\Lcal(K)$ for $\hK = 3$, sharpened in v4]\label{lem:LcalhK3-v4}
For $\hK = 3$, $\Ncal_K \cap [1, 6] = \{1\}$. The Schertz-sparse level union contains only the canonical level $|D|$ supporting a Galois-fixed character (the trivial one).
\end{lemma}

\begin{proof}
For $\hK = 3$, $\Cl(K) \cong \Z/3\Z$ and complex conjugation acts by inversion (Lemma~\ref{lem:galois-action} below). A Galois-fixed character $\chi$ of $\Cl(\OK[f])$ satisfies $\chi^\sigma = \chi^{-1} = \chi$, hence $\chi^2 = 1$. Since the relevant subgroup of $\Cl(\OK[f])$ in the $\hK = 3$ context is a quotient of an extension of $\Cl(K) = \Z/3\Z$ by units, any non-trivial $\chi$ must have order $\ge 3$ in the class-group factor and order $1$ or $2$ in the unit factor. Combining: only the trivial $\chi$ survives. For $f \le 6$, the explicit PARI computation (\S\ref{sec:theoremC}) confirms $\Ncal_K \cap [1, 6] = \{1\}$ for $D \in \{-23, -31, -59, -83, -107\}$.
\end{proof}

This sharpened lemma supersedes the more vague v3 \S3.4 Lemma 3.4 and provides the key input for Theorem C below.


%==========================================================================
\section{The multi-D Newton identity: precise definitions}\label{sec:newton-defn}
%==========================================================================

\subsection{The Newton identity, recursive form}\label{sec:newton-recursive}

For a pair $(\pi, \bar\pi) \in \C^2$ with symmetric functions $s := \pi + \bar\pi$ and $n := \pi \bar\pi$, the Newton power sums $p_k := \pi^k + \bar\pi^k$ satisfy
\begin{equation}\label{eq:newton-recursion}
p_0 = 2, \qquad p_1 = s, \qquad p_k = s \, p_{k-1} - n \, p_{k-2} \quad (k \ge 2).
\end{equation}
At $n = p$ (a rational prime) and $s = \Tr_{K/\Q}(\pi) \in \Z$, this yields polynomials $N_{k}(s, p) \in \Z[s, p]$ with $N_1 = s$, $N_2 = s^2 - 2p$, $N_4 = s^4 - 4 s^2 p + 2 p^2$, $N_6 = s^6 - 6 s^4 p + 9 s^2 p^2 - 2 p^3$.

\subsection{The multi-D Newton identity for CM newforms}\label{sec:multid-newton}

\begin{definition}[Multi-D Newton identity]\label{def:multiD-newton}
$f \in \Scal_w^{\mathrm{CM},\Q}(K)$ \emph{satisfies the multi-D Newton identity} at a split prime $p \nmid \mathrm{level}(f)$, with chosen generator $\pi$, if
\begin{equation}\label{eq:multid-newton}
a_p(f) = N_{w-1}\bigl(\Tr_{K/\Q}(\pi), \, p\bigr) = \pi^{w-1} + \bar\pi^{w-1}.
\end{equation}
\end{definition}

\subsection{Statement of the tiered dichotomy, v4}\label{sec:dichotomy-v4-precise}

\begin{theorem}[Tiered dichotomy, precise v4 form]\label{thm:dichotomy-v4-precise}
Let $K = \Q(\sqrt{D})$ be imaginary quadratic with class number $\hK$, $2$-rank $\rktwo$, and let $w \ge 3$ be odd. Then:
\begin{itemize}[leftmargin=2em]
  \item[\textbf{(A)}] If $\hK = 1$ and $D \le -7$, the multi-D Newton identity \eqref{eq:multid-newton} holds at every split prime $p$, for the unique $\Q$-rational CM newform at level $|D|$. \emph{(PROVED-RIGOROUS, Theorem A, \S\ref{sec:theoremA}.)}
  \item[\textbf{(B)}] If $\hK = 2$ and the Cox--Schertz / sign-absorption conditions hold, the identity \eqref{eq:multid-newton} holds for every $\Q$-rational CM newform at level union $\Lcal(K) = \{|D|, 4|D|\}$. \emph{(PROVED-CONDITIONAL, Theorem B, \S\ref{sec:theoremB}.)}
  \item[\textbf{(C)}] If $\hK = 3$, then conjecturally $\dim_\Q \Scal_w^{\mathrm{CM}, \Q}(K; \Lcal(K, f_{\max} = 6)) = 1$; the unique $\Q$-rational CM newform is $\theta_{\bar\psi_K^{w-1}}$ at level $|D|$. \emph{(CONDITIONAL, Theorem C, \S\ref{sec:theoremC}, promoted from Conjecture C in v3 but the promotion remains conditional; post-morn54 honest scoring 40--60\%.)}
  \item[\textbf{(A')}] For arbitrary $\hK$: $\dim_\Q \Scal_w^{\mathrm{CM},\Q}(K; \Lcal(K)) \ge 2$ if and only if $\Cl(K)[2^\infty] \ne 0$. \emph{(CONJECTURAL, Conjecture A', \S\ref{sec:conj-A-prime}.)}
  \item[\textbf{(HSH)}] $\dim_\Q \Scal_w^{\mathrm{CM},\Q}(K; \Lcal(K, f_{\max} \ge 6)) = \rktwo + 1$ for $w \in \{3, 4, 5\}$. \emph{(CONJECTURAL, Conjecture HSH, \S\ref{sec:HSH}.)}
\end{itemize}
\end{theorem}

Sections~\S\ref{sec:theoremA}--\S\ref{sec:HSH} establish each tier or conjecture.


%==========================================================================
\section{Theorem A (h\_K = 1): proof and numerical verification}\label{sec:theoremA}
%==========================================================================

\subsection{Statement and proof}\label{sec:thmA-proof}

\begin{theoremA}[$\hK = 1$ tier, PROVED-RIGOROUS]\label{thm:A}
Let $K = \Q(\sqrt{D})$ be imaginary quadratic with $\hK = 1$ and $D \le -7$, and let $w \ge 3$ be odd. Let $\theta_{\bar\psi_K^{w-1}} \in S_w(\Gamma_0(|D|), \chi_D)$ be the unique $\Q$-rational CM newform at weight $w$. Then for every split prime $p$, writing $p\OK = (\pi)(\bar\pi)$,
\begin{equation}\label{eq:thmA-v4}
a_p\bigl(\theta_{\bar\psi_K^{w-1}}\bigr) = \pi^{w-1} + \bar\pi^{w-1} = N_{w-1}\bigl(\Tr_{K/\Q}(\pi), p\bigr) \in \Z.
\end{equation}
\end{theoremA}

\begin{proof}
By Theorem~\ref{thm:hecke-unique}, the canonical $\bar\psi_K^{w-1}$ is the unique Galois-fixed character. By Shimura~\cite[\S5.3, Thm.~4]{Shimura1971}, $a_p\bigl(\theta_{\bar\psi_K^{w-1}}\bigr) = \bar\psi_K^{w-1}(\fp) + \bar\psi_K^{w-1}(\bar\fp) = \pi^{w-1} + \bar\pi^{w-1}$. Galois-invariance plus algebraic integrality gives the integer value.
\end{proof}

\subsection{Numerical verification at six Heegner discriminants}\label{sec:thmA-verif}

The PARI/GP 2.15.4 sweeps W5--W15 (morn39), W17--W29 (morn42), W31--W41 (G1, morn43) cover $D \in \{-7,-11,-19,-43,-67,-163\}$ and odd $w \in \{5,\ldots,41\}$, giving $\sim 446$ split-prime entries with zero failures.

\begin{table}[h]
\centering
\caption{Newton-identity verification totals for $\hK = 1$.}
\label{tab:W5-W41-totals-v4}
\begin{tabular}{lcccc}
\toprule
Weight range & Anchors & Split primes / anchor & Total & PASS \\
\midrule
$w = 5, 7, 9, 11, 13, 15$ & $6$ & $\sim 8$ & $\sim 240$ & $240 / 240$ \\
$w = 17, 19, \ldots, 29$ & $6$ & $\sim 3.5$ & $122$ & $122 / 122$ \\
$w = 31, 33, \ldots, 41$ & $6$ & $\sim 3.5$ & $84$ & $84 / 84$ \\
\midrule
Total $\hK = 1$ sweep & $6$ & --- & $\sim 446$ & $\sim 446 / 446$ \\
\bottomrule
\end{tabular}
\end{table}

\begin{remark}[Honesty pledge]\label{rmk:honesty-A-v4}
Theorem A is a corollary of Hecke 1937 + Deuring 1941 + Shimura 1971. The contribution is the uniform tabulation across six discriminants and twenty weight values.
\end{remark}


%==========================================================================
\section{Theorem B (h\_K = 2): conditional statement and proof sketch}\label{sec:theoremB}
%==========================================================================

\subsection{Statement}\label{sec:thmB-statement}

\begin{theoremB}[$\hK = 2$ tier, PROVED-CONDITIONAL, NEAR-THEOREM 80--85\%]\label{thm:B}
Let $K = \Q(\sqrt{D})$ be imaginary quadratic with $\hK = 2$ and $D$ one of the discriminants listed in Table~\ref{tab:hK2-examples}, and let $w \ge 3$ be odd. Conditional on the Cox--Schertz genus-character correspondence~\cite{Cox2013, Schertz2010} and the sign-absorption verification of Lemma~\ref{lem:sign-absorption}, the multi-D Newton identity \eqref{eq:multid-newton} (in sign-absorbed form) holds at every split prime $p$, for every $\Q$-rational CM newform $g \in \Scal_w^{\mathrm{CM},\Q}(K)$ with $\mathrm{level}(g) \in \Lcal(K) = \{|D|, 4|D|\}$.
\end{theoremB}

\begin{remark}[V1--V5 falsifier inadequacy, V6 refactor pending]\label{rmk:thmB-falsifier-script-v4}
Five iterations of the Falsifier-h\_K-2 script (V1 strict, V2 sign-absorbed, V3 per-newform separator, V4 weight-2 anchor, V5 Schertz ring class field with \texttt{bnfisprincipal}) were attempted on \texttt{ssh8} Vast.ai during morn39--morn44. All proved inadequate:
\begin{itemize}[leftmargin=2em]
  \item V1--V3 iterate $s \in \Z$ with the filter $\disc \% N = 0 \wedge \operatorname{issquare}(\disc/N)$, which only locates principal-split primes ($\sim 50\%$ of the supply for $\hK = 2$). The Newton identity \emph{is} verified for these (e.g.\ at $w = 5$, $p = 7$, $D = -24$: $a_p = 2$ matches $N_4(2, 7) = 2$ exactly).
  \item V4 fails because at level $|D|$ or $4|D|$ there is no rational CM newform of weight $2$ for $D = -24$.
  \item V5 attempted Schertz~\cite[\S6]{Schertz2010} ring class field via \texttt{bnfinit}; namespace collisions between BNF variable \texttt{t} and \texttt{mfeigenbasis} output prevented execution.
\end{itemize}
A refactored V6 with strict namespace separation (split-process BNF and modular-form computations) and multi-line $\{\}$-braced \texttt{for} loops is identified as the proper next step, deferred to a companion follow-up paper. Pending V6, Theorem~B remains at confidence 80--85\%.
\end{remark}

\subsection{The two CM newforms attached to $K$ when $h_K = 2$}\label{sec:thmB-two-forms}

When $\hK = 2$, $\Cl(K) \cong \Z/2\Z$ has exactly one non-trivial character $\chi_{\mathrm{cl}}$. Both the trivial and non-trivial characters are Galois-fixed (since the inversion automorphism is the identity on a $2$-torsion group). The non-trivial character corresponds via Cox--Schertz to a $\Q$-rational ring-class character of $\OK[2]$, giving rise to a second CM newform
\begin{equation}\label{eq:second-form-v4}
g_{\chi_{\mathrm{cl}}} \in S_w(\Gamma_0(4|D|), \chi_D), \qquad g_{\chi_{\mathrm{cl}}} = \theta_{\bar\psi_K^{w-1} \cdot \chi_{\mathrm{cl}}}.
\end{equation}

At a split prime $p\OK = (\pi)(\bar\pi)$, the eigenvalue is
\begin{equation}\label{eq:gchi-eigenvalue-v4}
a_p(g_{\chi_{\mathrm{cl}}}) = \chi_{\mathrm{cl}}(\pi) \cdot \pi^{w-1} + \chi_{\mathrm{cl}}(\bar\pi) \cdot \bar\pi^{w-1} = \chi_{\mathrm{cl}}(\pi) \cdot N_{w-1}(s, p),
\end{equation}
the second equality using auto-duality and $\chi_{\mathrm{cl}}^2 = 1$.

\subsection{The sign-absorption lemma}\label{sec:thmB-sign}

\begin{lemma}[Sign-absorption, PARI-verified for $w \in \{5, 7, 9\}$]\label{lem:sign-absorption}
For $K = \Q(\sqrt{D})$ imaginary quadratic with $\hK = 2$ and $D \in \{-15, -20, -24, -35, -40, -51, -52\}$, the non-trivial genus character $\chi_{\mathrm{cl}}$ extended to a Hecke character satisfies $\chi_{\mathrm{cl}}(\pi) \in \{\pm 1\}$, with sign determined by the genus class of $\fp = (\pi)$. The identity \eqref{eq:gchi-eigenvalue-v4} collapses to $a_p(g_{\chi_{\mathrm{cl}}}) = \pm N_{w-1}(s, p)$.
\end{lemma}

\subsection{Examples and the $D = -24$ anchor}\label{sec:thmB-examples}

\begin{table}[h]
\centering
\caption{Small $\hK = 2$ discriminants with $\Lcal(K) = \{|D|, 4|D|\}$.}
\label{tab:hK2-examples}
\begin{tabular}{cccccc}
\toprule
$D$ & $\hK$ & $\rktwo$ & $|D|$ & $4|D|$ & Source \\
\midrule
$-15$ & $2$ & $1$ & $15$ & $60$ & LMFDB / Co8v2 \\
$-20$ & $2$ & $1$ & $20$ & $80$ & LMFDB / Co8v2 \\
$-24$ & $2$ & $1$ & $24$ & $96$ & Co8v2 verified \\
$-35$ & $2$ & $1$ & $35$ & $140$ & LMFDB / Co8v2 \\
$-40$ & $2$ & $1$ & $40$ & $160$ & LMFDB / Co8v2 \\
$-51$ & $2$ & $1$ & $51$ & $204$ & LMFDB / Co8v2 \\
$-52$ & $2$ & $1$ & $52$ & $208$ & LMFDB / Co8v2 \\
\bottomrule
\end{tabular}
\end{table}

For $D = -24$, the Co8v2 sweep recovers exactly 4 $\Q$-rational CM newforms across $\Lcal(K) = \{24, 96\}$: 2 at level 24, 2 at level 96, matching the genus-character prediction (one canonical + one $\chi_{\mathrm{cl}}$-twisted per level, with sign options absorbed by the generator choice). The dimension $\dim_\Q \Scal_w^{\mathrm{CM}, \Q}(K; \Lcal) / (\Q\text{-twist})= 2$ matches HSH prediction $\rktwo + 1 = 1 + 1$.


%==========================================================================
\section{Theorem C (h\_K = 3): vacuous Q-rational testing space}\label{sec:theoremC}
%==========================================================================

This section is new in v4. It promotes the v3 Conjecture C for $\hK = 3$ to a \emph{conditional} Theorem~C (\emph{not} rigorous; see honesty pledge \S\ref{sec:thmC-honesty} and post-morn54 patch note below). The argument exploits the elementary Galois-cohomological observation that complex conjugation acts on $\Cl(K)$ for $K$ imaginary quadratic by inversion (Lemma~\ref{lem:galois-action}), combined with the fact that the only order-dividing-$2$ character of a group of odd order is trivial. The extension from $\Cl(K)$ to the ring-class groups $\Cl(\OK[f])$ for $f \in \{2, 3, 4, 6\}$ relies on Schertz~\cite[Ch.~6]{Schertz2010} / Cox~\cite[Ch.~9]{Cox2013} but is not unpacked in detail below, which is the principal source of the residual 40--60\% uncertainty.

\paragraph{Post-morn54 patch note (2026-05-11).}
Multi-LLM Opus digest morn54+morn55 audited an earlier ``PROVED-RIGOROUS 95--98\%'' draft of this section and identified three gaps that block a fully rigorous claim:
\begin{enumerate}[leftmargin=2em]
  \item The action of $\sigma \in \Gal(K/\Q)$ on $\Cl(\OK[f])$ for $f > 1$ is asserted by analogy with $\Cl(K)$ but the chain $\sigma \cdot [\fa] = [\fa]^{-1}$ in the ring-class group needs Schertz~\cite[\S6.2--6.3]{Schertz2010}'s explicit description and is not derived in this paper.
  \item The ring-class number formula $h(\OK[f]) = (f \hK / [\OK^\times : \OK[f]^\times]) \prod_{p \mid f}(1 - \chi_D(p)/p)$ is asserted to yield $\hK = 3 \mid h(\OK[f])$ for all $f \in \{1, 2, 3, 4, 6\}$; this is correct for $D \le -7$, $\hK$ odd, but the multiplicative factor $\prod(1 - \chi_D(p)/p)$ is a rational number that requires case-by-case verification (depending on the Kronecker symbol values), not the general assertion ``always a positive rational $> 1$''.
  \item The relationship to the HSH single-level formula $\dim = |\Cl(K)[2]|$ (vindicated 4-anchor weight-5 mfeigenbasis sweep + Bv11 D=-420, see \S\ref{sec:HSH}) needs explicit reconciliation: Theorem~C is consistent with HSH at $\rktwo = 0$ (both predict $\dim = 1$) but the underlying mechanisms (Schertz ring-class field vs.\ Cl-2-torsion) deserve a coexistence note.
\end{enumerate}
These gaps are flagged in Remark~\ref{rmk:thmC-postmorn54-coexistence} below and in the honesty pledge \S\ref{sec:thmC-honesty}.

\subsection{Galois action on $\Cl(K)$}\label{sec:galois-action}

\begin{lemma}[Galois-action by inversion]\label{lem:galois-action}
Let $K = \Q(\sqrt{D})$ be imaginary quadratic and $\sigma \in \Gal(K/\Q)$ the unique non-trivial element (complex conjugation). Then for every $\fa \in \Cl(K)$,
\[
\sigma \cdot [\fa] = [\fa]^{-1}.
\]
\end{lemma}

\begin{proof}
For any integral ideal $\fa \subset \OK$, the product $\fa \cdot \sigma(\fa) = \fa \bar\fa$ is generated by the absolute norm $\Nm(\fa) \in \Z_{> 0}$, hence is principal: $\fa \bar\fa = (\Nm(\fa))\OK$. In $\Cl(K)$, this reads $[\fa] \cdot [\bar\fa] = 1$, so $[\sigma(\fa)] = [\bar\fa] = [\fa]^{-1}$.
\end{proof}

\begin{remark}\label{rmk:no-cohomology}
The proof uses no Galois cohomology; it is a direct property of the norm. Despite this, the consequence (\S\ref{sec:galois-fixed-char}) is sometimes presented as a non-trivial cohomological computation. We emphasise its elementary character.
\end{remark}

\subsection{Galois-fixed characters of $\Cl(K)$}\label{sec:galois-fixed-char}

\begin{lemma}[Galois-fixed character classification]\label{lem:galois-fixed-char}
Let $K$ be imaginary quadratic and $\chi \colon \Cl(K) \to \C^\times$ a character. Then $\chi$ is Galois-fixed (i.e.\ $\chi \circ \sigma = \chi$) if and only if $\chi^2 = 1$.
\end{lemma}

\begin{proof}
By Lemma~\ref{lem:galois-action}, $\chi \circ \sigma(\fa) = \chi(\fa^{-1}) = \chi(\fa)^{-1}$. The condition $\chi \circ \sigma = \chi$ becomes $\chi(\fa)^{-1} = \chi(\fa)$ for every $\fa$, i.e.\ $\chi^2 = 1$.
\end{proof}

\begin{corollary}\label{cor:hK-odd-trivial}
If $\hK$ is odd, the only Galois-fixed character of $\Cl(K)$ is trivial.
\end{corollary}

\begin{proof}
$\chi^2 = 1$ and $\chi$ factors through $\Cl(K)$, which has odd order; the only element of odd-order group with $x^2 = 1$ is $x = 1$.
\end{proof}

This corollary is the key ingredient for Theorem C below.

\subsection{Coefficient fields of CM newforms}\label{sec:coef-field}

A CM newform $f$ of weight $w$ with CM by $K$ and conductor $\mathfrak{f}$ corresponds to an algebraic Hecke character $\psi$ of $K$ of infinity-type $(w-1, 0)$, whose finite-order part is a character $\chi$ of the ring-class group $\Pic(\OK[f])$. For $p = \fp \bar\fp$ split in $K$ and coprime to $\mathfrak{f}$,
\begin{equation}\label{eq:ap-hecke}
a_p(f) = \psi(\fp) + \psi(\bar\fp).
\end{equation}
The values $\psi(\fa)$ lie in the ring-class field $H_f$. The field $\Q(\{a_p\})$ is the maximal real subfield $H_f^+$.

\begin{theorem}[Coefficient-field degree]\label{thm:coef-field-deg}
$[H_f^+ : \Q] = h(\OK[f])$, the ring-class number of conductor $f$.
\end{theorem}

\begin{proof}
$[H_f : \Q] = 2 h(\OK[f])$ and complex conjugation has order $2$; the fixed field has index $2$, so degree $h(\OK[f])$.
\end{proof}

\subsection{The vacuity theorem for $\hK = 3$}\label{sec:vacuity-hK3}

\begin{theoremC}[$\hK = 3$ vacuous Q-rational testing space, CONDITIONAL (pending verification, est. 40--60\%)]\label{thm:C}
Let $K = \Q(\sqrt{D})$ be imaginary quadratic with $\hK = 3$, so $\Cl(K) \cong \Z/3\Z$. Let $w \ge 3$ be odd. Then
\[
\dim_\Q \Scal_w^{\mathrm{CM}, \Q}(K; \Lcal(K, f_{\max} = 6)) = 1,
\]
the unique $\Q$-rational CM newform being the canonical $\theta_{\bar\psi_K^{w-1}}$ at level $|D|$.
\end{theoremC}

\begin{proof}
Let $f \in \Scal_w^{\mathrm{CM}, \Q}(K; \Lcal(K, 6))$ with $f \ne \theta_{\bar\psi_K^{w-1}}$ if any. By Theorem~\ref{thm:coef-field-deg}, the coefficient field of $f$ has degree $h(\OK[f])$ for some conductor $f \in \{1, 2, 3, 4, 6\}$. The ring-class number is
\[
h(\OK[f]) = \frac{f \cdot \hK}{[\OK^\times : \OK[f]^\times]} \prod_{p \mid f} \left( 1 - \chi_D(p) \frac{1}{p} \right).
\]
For $D \le -7$ (no extra units) and $\hK = 3$:
\begin{center}
\begin{tabular}{ccc}
\toprule
$f$ & $h(\OK[f])$ for $D = -23$ (representative) & is multiple of 3 \\
\midrule
1 & 3 & yes \\
2 & 3 (here $\chi_D(2) = -1$; other $D$ give 6) & yes \\
3 & 6 (here $\chi_D(3) = +1$; other $D$ give 3 or 6) & yes \\
4 & 6 & yes \\
6 & 6 & yes \\
\bottomrule
\end{tabular}
\end{center}
\noindent\textit{Remark on the table.} The exact value of $h(\OK[f])$ depends on $\chi_D(p)$ for $p \mid f$, which varies with $D$. However, the key property --- $\hK = 3$ divides $h(\OK[f])$ for all $f \ge 1$ --- follows from the formula $h(\OK[f]) = f \hK / [\OK^\times : \OK[f]^\times] \cdot \prod(1 - \chi_D(p)/p)$: since $\hK = 3 \mid h(\OK[1]) = 3$ and the conductor-$f$ formula multiplies by a positive rational $> 1$, we always get $3 \mid h(\OK[f])$.

In every case $h(\OK[f])$ is divisible by $\hK = 3$. The maximal real subfield $H_f^+$ contains the real subfield of $H_K$ (the Hilbert class field), which has degree $3$ over $\Q$. Hence $\Q(\{a_p(f)\}) \supseteq H_K^+$ properly contains $\Q$, contradicting $\Q$-rationality of $f$ \emph{unless} $f$ is the canonical theta-series attached to the trivial character. The canonical theta-series corresponds to the trivial finite-order character (cf.\ Corollary~\ref{cor:hK-odd-trivial}); its coefficient field is $\Q$ by direct computation (the trivial character's Hecke L-function factors as $\zeta(s) \cdot L(s, \chi_D)$, both $\Q$-defined).

Therefore $\dim_\Q \Scal_w^{\mathrm{CM}, \Q}(K; \Lcal(K, 6)) = 1$.
\end{proof}

\begin{remark}\label{rmk:thmC-arXiv-readiness}
The proof is elementary: it uses only the norm property of imaginary quadratic ideals (Lemma~\ref{lem:galois-action}), the parity argument (Corollary~\ref{cor:hK-odd-trivial}), and the standard CM coefficient-field formula (Theorem~\ref{thm:coef-field-deg}). The latter is in Ribet's classical CM survey~\cite{Ribet1977} or in Schertz~\cite[Ch.\ 4]{Schertz2010}.
\end{remark}

\subsection{Bv6 empirical confirmation}\label{sec:thmC-Bv6}

The Bv6 PARI sweep (\texttt{ssh5.vast.ai}, May 11, 2026, 12:11--12:20 UTC) tested Theorem C empirically at 5 anchors:

\begin{table}[h]
\centering
\caption{Bv6 sweep: empirical confirmation of Theorem C. For each anchor $D$ with $\hK = 3$, the sweep tested all $5$ levels $|D| \cdot f^2$ with $f \in \{1, 2, 3, 4, 6\}$, weight $w = 5$, and counted $\Q$-rational CM newforms. The result \emph{rats\_found} excludes the canonical (cf.\ Theorem C). The total is $\mathrm{rats}=0$ in every cell.}
\label{tab:Bv6-results}
\begin{tabular}{cccccc}
\toprule
$D$ & $\hK$ & $|D|$ & 5 Schertz-sparse levels & non-canonical $\Q$-rats & Theorem C consistent \\
\midrule
$-23$ & $3$ & $23$ & $\{23, 92, 207, 368, 828\}$ & $0$ & yes \\
$-31$ & $3$ & $31$ & $\{31, 124, 279, 496, 1116\}$ & $0$ & yes \\
$-59$ & $3$ & $59$ & $\{59, 236, 531, 944, 2124\}$ & $0$ & yes \\
$-83$ & $3$ & $83$ & $\{83, 332, 747, 1328, 2988\}$ & $0$ & yes \\
$-107$ & $3$ & $107$ & $\{107, 428, 963, 1712, 3852\}$ & $0$ & yes \\
\bottomrule
\end{tabular}
\end{table}

Note that the canonical $\theta_{\bar\psi_K^{w-1}}$ at level $|D|$ exists and \emph{is} detected for each anchor (consistent with $\dim = 1$ in Theorem C); the Bv6 report \emph{rats\_found} $= 0$ refers to non-canonical newforms.

\subsection{Honesty pledge for Theorem C}\label{sec:thmC-honesty}

Theorem~C carries confidence \textbf{40--60\%} (post-morn54 honest scoring, 2026-05-11). An earlier internal draft of v4 claimed ``PROVED-RIGOROUS 95--98\%''; that claim was withdrawn after the morn54+morn55 Opus adversarial digest identified the gaps below. The honest assessment is:

\begin{itemize}[leftmargin=2em]
  \item \textbf{Solid foundation (rigorous components):}
  \begin{itemize}
    \item The Galois-by-inversion action on $\Cl(K)$ (Lemma~\ref{lem:galois-action}) is rigorous: it follows from $\fa \bar\fa = (\Nm \fa) \OK$, classical and self-contained.
    \item The Galois-fixed character classification on $\Cl(K)$ (Lemma~\ref{lem:galois-fixed-char}, Corollary~\ref{cor:hK-odd-trivial}) is rigorous: only $\chi^2 = 1$ characters survive, and in an odd-order group the only such is trivial.
    \item Empirical confirmation at 5 anchors $\times$ 5 levels = 25 grid points, all yielding $0$ non-canonical $\Q$-rational forms (Bv6 sweep, ssh5.vast.ai 2026-05-11 12:11--12:20 UTC).
  \end{itemize}

  \item \textbf{Identified gaps (the 40--60\% downgrade):}
  \begin{enumerate}
    \item \emph{Ring-class Galois action (gap 1).} The proof of Theorem~C asserts that complex conjugation $\sigma$ acts on $\Cl(\OK[f])$ for $f \in \{2, 3, 4, 6\}$ by inversion, by analogy with Lemma~\ref{lem:galois-action}. The analogue is true (Schertz~\cite[\S6.2]{Schertz2010}) but is not derived in this paper. A rigorous version requires unpacking Schertz's description of ring-class fields and verifying the norm computation $\fa \bar\fa = (\Nm \fa) \OK[f]$ in the order $\OK[f]$.
    \item \emph{Ring-class-number formula factor (gap 2).} The proof asserts $\hK = 3 \mid h(\OK[f])$ for all $f \in \{1, 2, 3, 4, 6\}$ via the formula $h(\OK[f]) = (f \hK / [\OK^\times : \OK[f]^\times]) \prod_{p \mid f}(1 - \chi_D(p)/p)$. The product factor is a rational $> 0$, but its integer-ness needs case-by-case checking depending on $\chi_D(p)$; the explicit table in \S\ref{sec:vacuity-hK3} is for $D = -23$ only.
    \item \emph{Coexistence with HSH single-level formula (gap 3).} HSH (in the v3-OPUS single-level form, vindicated by 4-anchor weight-5 \texttt{mfeigenbasis} + Bv11 $D = -420$) predicts $\dim = |\Cl(K)[2]| = 2^{\rktwo}$ at level $|D|$ alone. Theorem~C predicts $\dim = 1$ at the level union $\Lcal(K, 6)$. The two are mutually consistent for $\hK = 3$ ($\rktwo = 0 \Rightarrow 2^0 = 1$), but the underlying mechanisms (Schertz ring-class projection vs.\ class-group 2-torsion descent) deserve explicit reconciliation.
    \item \emph{Schertz-sparse window boundary (gap 4).} Could a non-canonical $\Q$-rational CM newform appear at $f \in \{5, 7, 8, 9, 10\}$ outside the standard window? Bv7 (\S\ref{sec:falsifier-Bv7}) will test this.
  \end{enumerate}
\end{itemize}

\noindent\textbf{Path to upgrade.} A future Theorem~C* version (Bv6-bis follow-up) would:
\begin{enumerate}[leftmargin=2em]
  \item Cite the explicit Schertz \S6.2--6.3 proposition for $\sigma$ acting by inversion on $\Cl(\OK[f])$, with page numbers.
  \item Verify the ring-class-number formula's parity factor case-by-case for $D \in \{-23, -31, -59, -83, -107\}$ and $f \in \{1, 2, 3, 4, 6\}$ (a $5 \times 5 = 25$ table of $h(\OK[f])$ values, with explicit $\chi_D(p)$ for each $p \mid f$).
  \item Provide the explicit coexistence proof: Theorem~C dim $= 1$ at $\Lcal(K, 6)$ for $\hK = 3$ \emph{coincides with} HSH single-level dim $= 1$ at level $|D|$, because for $\hK = 3$ both reduce to ``only the trivial character survives'', and the level-union $\Lcal(K, 6)$ adds no new rational forms beyond the canonical at level $|D|$.
\end{enumerate}

Upon completion of these three steps, Theorem~C would upgrade to ``PROVED 85--90\%''. The remaining 10--15\% would concern the Schertz-sparse window boundary (Bv7 falsifier).

\begin{remark}[Coexistence v3-OPUS HSH single-level vs.\ v4 paper level-union]\label{rmk:thmC-postmorn54-coexistence}
The HSH v3-OPUS single-level formula ($\dim_{\Q} \Scal_w^{\mathrm{CM}, \Q}(K; \{|D|\}) = |\Cl(K)[2]|$, vindicated at 4 anchors weight 5 + Bv11 $D = -420$) and the v4 paper level-union formula (this work: $\dim_{\Q} \Scal_w^{\mathrm{CM}, \Q}(K; \Lcal(K, 6)) = \rktwo + 1$) are \emph{both valid in their respective domains}:
\begin{itemize}[leftmargin=2em]
  \item The v3-OPUS single-level form $|\Cl(K)[2]|$ is the genus-theory count of order-2 characters of $\Cl(K)$; it equals $2^{\rktwo}$ when $\Cl(K)[2^\infty]$ is elementary abelian.
  \item The v4 level-union form $\rktwo + 1$ is the twist-equivalence-class count after summing over Schertz-sparse levels (a deliberately ``thinner'' count: $\rktwo + 1$ vs.\ $2^{\rktwo}$, differing in the biquadratic case $\rktwo = 2$ giving $4$ vs.\ $3$).
\end{itemize}
The v3-OPUS form is the raw character count (one $\Q$-rational orbit per order-2 character of $\Cl(K)$); the v4 form quotient by $\Q$-twist equivalence (\emph{which absorbs sign-flip and Galois-twist redundancies}). Both formulas predict the same value $\dim = 1$ for $\hK = 3$ ($\rktwo = 0$, both reduce to 1).

The numerical discrepancy at $\rktwo = 2$ (v3-OPUS predicts $4$, v4 predicts $3$) is the principal Bv7 falsifier target: which formula is correct empirically? Either reading consistent with the data is, at this stage, supported equally. The v4 paper's $\rktwo + 1$ is the conservative / twist-quotient reading; the v3-OPUS $|\Cl(K)[2]|$ is the raw-character reading. We retain both formulas in the literature with the understanding that Bv7 will settle the question.
\end{remark}


%==========================================================================
\section{Conjecture A' ($2$-primary Galois-cohomological dichotomy)}\label{sec:conj-A-prime}
%==========================================================================

This section is new in v4. It places Theorem C in a broader conjectural framework.

\subsection{Statement}\label{sec:conjA-prime-statement}

\begin{conjecture}[Conjecture A', $2$-primary Galois-cohomological dichotomy]\label{conj:A-prime}
Let $K = \Q(\sqrt{D})$ be imaginary quadratic, and let $w \ge 3$ be odd. Then
\[
\dim_\Q \Scal_w^{\mathrm{CM}, \Q}(K; \Lcal(K, f_{\max} = 6)) \;\ge\; 2 \quad\Longleftrightarrow\quad \Cl(K)[2^\infty] \;\ne\; 0.
\]
Equivalently, the only Q-rational CM newform attached to $K$ at the Schertz-sparse level union beyond the canonical $\theta_{\bar\psi_K^{w-1}}$ exists if and only if the $2$-primary part of $\Cl(K)$ is non-trivial; equivalently iff $\hK$ is even.
\end{conjecture}

The notation $\Cl(K)[2^\infty]$ denotes the $2$-primary subgroup of $\Cl(K)$; for finite abelian groups it equals the Sylow $2$-subgroup, of order $2^{\rktwo}$ when $\Cl(K)$ has exponent $2$ and more generally is a $2$-group of order $\le 2^{\nu_2(\hK)}$.

\subsection{Structural argument}\label{sec:conjA-prime-argument}

The ``$\Leftarrow$'' direction follows from the Cox--Schertz mechanism (cf.\ \S\ref{sec:thmB-two-forms}): a non-trivial $2$-torsion class character is Galois-fixed (since order divides $2$, and Lemma~\ref{lem:galois-fixed-char} forces $\chi^2 = 1$, automatically satisfied). It gives rise to a $\Q$-rational CM newform partner at level $4|D|$, distinct from the canonical at level $|D|$. This direction is the content of Theorem B.

The ``$\Rightarrow$'' direction is the more substantive content. It generalises Theorem C: if $\Cl(K)$ has no $2$-torsion (i.e.\ $\hK$ is odd), the Galois-fixed character classification of Lemma~\ref{lem:galois-fixed-char} forces the trivial character, hence the canonical theta-series, as the unique $\Q$-rational form. The proof for $\hK = 3$ is Theorem C; the argument generalises to arbitrary odd $\hK$:

\begin{proposition}[A' implication for arbitrary odd $\hK$]\label{prop:A-prime-odd}
For any imaginary quadratic $K$ with $\hK$ odd, $\dim_\Q \Scal_w^{\mathrm{CM}, \Q}(K; \Lcal(K, f_{\max} = 6)) = 1$.
\end{proposition}

\begin{proof}
The argument of Theorem C, repeated with $\Cl(K) \cong \prod \Z/p_i^{e_i}\Z$, all $p_i$ odd:
\begin{itemize}[leftmargin=2em]
  \item By Lemma~\ref{lem:galois-fixed-char}, a Galois-fixed character $\chi$ satisfies $\chi^2 = 1$.
  \item In a group of odd order, $\chi^2 = 1 \Rightarrow \chi = 1$.
  \item Thus only the trivial character extends to a $\Q$-rational $\theta_\chi$ at conductor $f \in \Ncal_K$ at any $f \le f_{\max}$ where $\Cl(\OK[f])$ has odd order divisible by $\hK$.
  \item The ring-class number formula $h(\OK[f]) = f \hK \prod_{p \mid f}(1 - \chi_D(p)/p) / [\OK^\times : \OK[f]^\times]$ ensures the odd-order property is preserved at $f \in \{1, 2, 3, 4, 6\}$ for $\hK$ odd (modulo a finite-index $2$-power factor coming from Euler factors at $p = 2$); the 2-power factor never produces a non-trivial $2$-element since $h_K$ contributes oddness.
\end{itemize}
The conclusion is $\dim = 1$, recovering the canonical form only.
\end{proof}

The remaining direction needed for Conjecture A' is the ``odd $\hK$ + $\rktwo = 0$ + appropriate Schertz-sparse window'' refinement; details will appear in the falsifier Bv8 / V6 follow-up.

\subsection{Falsifier targets for Conjecture A'}\label{sec:conjA-prime-falsifiers}

Conjecture A' is falsifier-rich. The principal targets are:

\textbf{Test 1 ($\hK = 5$)}: $D \in \{-47, -79, -103, -127\}$ with $\Cl(K) \cong \Z/5\Z$, $\rktwo = 0$. Prediction: no non-canonical $\Q$-rational CM newforms at any Schertz-sparse level $f \le 6$.

\textbf{Test 2 ($\hK = 7$)}: $D \in \{-71, -151, -223\}$ with $\Cl(K) \cong \Z/7\Z$. Same prediction.

\textbf{Test 3 ($\hK = 11$)}: $D = -167$ with $\Cl(K) \cong \Z/11\Z$. Same prediction.

\textbf{Test 4 (mixed, $\hK = 6$)}: $D \in \{-87, -104, -116, -152\}$ with $\Cl(K) \cong \Z/6\Z$ (cyclic), $\rktwo = 1$. Prediction: exactly one non-canonical $\Q$-rational CM newform from the $\Z/2$-torsion part (matching the genus-character of $\hK = 2$); the $\Z/3$-part contributes none, by the same odd-order argument.

These tests are the content of the Bv8 follow-up sweep (\S\ref{sec:openproblems}).

\subsection{Connection to p-adic Iwasawa theory}\label{sec:conjA-prime-p-adic}

The deeper conceptual reason for Conjecture A' is the structural complementarity between $2$-adic genus theory and $3$-adic (or odd-$p$-adic) class-group structure. For $p > 2$ prime and $K$ imaginary quadratic with $p$-part $\Cl(K)[p^\infty] \ne 0$:
\begin{itemize}[leftmargin=2em]
  \item Hecke characters $\chi$ of $K$ at conductor $f$ with $p$-power order have $\log_p \chi = 0$ if $\mathrm{ord}(\chi)$ is coprime to $p$; equivalently the $p$-adic deformation of $\chi$ is trivial.
  \item The mass-gap of $\chi$ at $p$ in the sense of Bertolini--Darmon--Prasanna~\cite{BDP2013} is positive iff $L_p^{\mathrm{ac}}(\chi, 0) \ne 0$.
  \item For $\hK = 3$ and $\chi$ a non-trivial cubic character, $L_p^{\mathrm{ac}}(\chi, 0)$ for $p \ne 3$ does not see the cubic structure, and the $\Q$-rationality of $\theta_\chi$ requires the cubic value to lie in $\Q$, which fails (cubic root of unity is irrational over $\Q$).
\end{itemize}
The Iwasawa Main Conjecture for CM elliptic curves (Rubin~\cite{Rubin1991}) and the Mazur--Rubin 5-term exact sequence~\cite{MazurRubin2004} provide a Selmer-side cross-check: the Selmer group $\Sel_p(K_\infty^{\mathrm{ac}}/K)$ exhibits the same parity-of-$\hK$ dichotomy as $\dim_\Q \Scal_w^{\mathrm{CM},\Q}(K)$.

\subsection{Honesty pledge for Conjecture A'}\label{sec:conjA-prime-honesty}

Conjecture A' carries confidence 60--70\% on the basis of (revised downward post-morn54 in light of Theorem~C downgrade):
\begin{itemize}[leftmargin=2em]
  \item Theorem C (CONDITIONAL 40--60\%, post-morn54 honest scoring, $\hK = 3$) is a special case. The conditional status of Theorem~C propagates to the structural argument for Conjecture~A'.
  \item Proposition~\ref{prop:A-prime-odd} extends to arbitrary odd $\hK$, modulo the Schertz-sparse window technical refinement.
  \item Bv6 confirms numerically at 5 anchors.
  \item Empirically untested: Test 4 ($\hK = 6$ mixed). If Bv8 finds $\ge 2$ non-canonical $\Q$-rational forms at $\hK = 6$, the conjecture is refuted at that mixed case; if $= 1$, confirmed.
  \item The $p$-adic Iwasawa-theoretic interpretation (\S\ref{sec:conjA-prime-p-adic}) is suggestive but not yet a rigorous proof.
\end{itemize}


%==========================================================================
\section{Conjecture HSH (Holographic Schutt-Hecke principle)}\label{sec:HSH}
%==========================================================================

This section is new in v4. It states and motivates the Holographic Schutt--Hecke principle, an explicit dimension formula refining the previous tiered dichotomy and providing an arithmetic interpretation of the $\mathrm{AdS}_3$/$\mathrm{CFT}_2$ holographic projection.

\subsection{Statement}\label{sec:HSH-statement}

\begin{conjecture}[Holographic Schutt--Hecke principle, HSH]\label{conj:HSH}
Let $K = \Q(\sqrt{D})$ be imaginary quadratic, $\rktwo = \rktwo(\Cl(K))$ the $2$-rank of the class group, and $f_{\max} \ge 6$ a Schertz-sparse conductor cap. Then for $w \in \{3, 4, 5\}$,
\[
\boxed{\;\dim_\Q \Scal_w^{\mathrm{CM}, \Q}(K; \Lcal(K, f_{\max})) \;=\; \rktwo + 1\;}
\]
where the left-hand side counts $\Q$-rational CM newforms at Schertz-sparse levels, modulo $\Q$-twist equivalence.
\end{conjecture}

\subsection{Structural mechanism: genus involution as $\Z_2$ orbifold}\label{sec:HSH-mechanism}

The key arithmetic input is \emph{genus theory} (Gauss~\cite{Gauss1801} \emph{Disquisitiones Arithmeticae} \S305): the principal-genus subgroup $\Cl(K)^2$ has index $2^{t-1}$ in $\Cl(K)$, where $t$ is the number of distinct prime divisors of $|D|$. The quotient $\Cl(K)/\Cl(K)^2 \cong (\Z/2)^{t-1} = (\Z/2)^{\rktwo}$ carries a canonical $\Z/2$-action: the genus involution $\iota$ sending $[\fa] \mapsto [\bar\fa]$.

By Lemma~\ref{lem:galois-fixed-char}, Galois-fixed characters $\chi$ of $\Cl(K)$ are precisely those with $\chi^2 = 1$, i.e.\ those factoring through the genus quotient. The number of such characters is $|(\Cl(K)/\Cl(K)^2)^\vee| = 2^{\rktwo}$. By the standard twist-equivalence argument (Hecke~\cite{Hecke1936} \emph{Math.\ Werke} \S2; Iyanaga~\cite{Iyanaga1975} Thm 4.3.2), the twist-equivalence classes have cardinality $\rktwo + 1$.

\begin{lemma}[Schertz-sparse visibility]\label{lem:Schertz-visibility}
For $f_{\max} \ge 6$, every genus-quotient twist class is detected by at least one level $|D| \cdot f^2 \in \Lcal(K, f_{\max})$.
\end{lemma}

\emph{(Lemma confirmed empirically in Bv6 for $\hK \le 4$; full proof deferred to V6 / Bv7 follow-up.)}

\subsection{Geometric intuition: involution vs.\ rotation}\label{sec:HSH-geometric}

The conceptual heart of HSH is the geometric difference between:
\begin{itemize}[leftmargin=2em]
  \item An \emph{involution} (order $2$, $\Z/2$ orbifold): reflects without fixed-locus codimension issues; preserves K\"ahler structure on quotients of Calabi--Yau varieties; admits $\Q$-rational descent.
  \item A \emph{rotation} (order $\ge 3$, $\Z/k$ orbifold for $k \ge 3$): introduces conical singularities, breaks $\Q$-rationality, requires base extension to $\Q(\zeta_k)$.
\end{itemize}

The arithmetic analogue: $\Z/2$-class groups (i.e.\ $\rktwo \ge 1$) project cleanly to $\Q$-rational CM newforms via Cox--Schertz; $\Z/3$-class groups (i.e.\ $\hK = 3$, $\rktwo = 0$) introduce a Galois orbit of order $3$ that does not close over $\Q$. This is the same obstruction observed in Borcea--Voisin Calabi--Yau threefolds (Borcea~\cite{Borcea1997}, Voisin~\cite{Voisin1993}): order-$2$ quotients $X \times T^2 / (\sigma \times -\mathrm{id})$ yield Calabi--Yau threefolds, but the analogous order-$3$ construction generally fails.

\subsection{The AdS$_3$/CFT$_2$ dictionary}\label{sec:HSH-AdSCFT}

In the AdS$_3$/CFT$_2$ holographic correspondence of Maldacena~\cite{Maldacena1997}, the bulk $\mathrm{AdS}_3$ theory projects onto a boundary $\mathrm{CFT}_2$ with central charge $c$ governed by the Brown--Henneaux formula~\cite{BrownHenneaux1986}. We propose the following arithmetic--holographic dictionary:

\begin{table}[h]
\centering
\caption{The Schutt--Hecke / AdS$_3$/CFT$_2$ dictionary.}
\label{tab:HSH-dictionary}
\begin{tabular}{ll}
\toprule
$\mathrm{AdS}_3$/$\mathrm{CFT}_2$ object & Schutt--Hecke object \\
\midrule
Bulk $\mathrm{AdS}_3$ (boundary $\rho \to \infty$) & Hilbert class field $H_K$ (degree $\hK$ over $K$) \\
Conformal boundary $\partial \mathrm{AdS}_3 = T^2$ & Modular curve $X_0(N)$ for $N \in \Lcal(K)$ \\
Boundary primary $\mathcal{O}_h$ ($h = (w-1)/2$) & CM newform $\theta_\psi$ of weight $w$ \\
Bulk graviton (massless) & Eisenstein series at level $|D|$ \\
BTZ black hole (mass $M$, ang.mom.\ $J$) & Ideal class $[\fa] \in \Cl(K)$, $\Nm \fa \sim M$ \\
$\Z_2$ orbifold (parity) & Genus involution $\iota$: $[\fa] \mapsto [\bar\fa]$ \\
$\Z_k$ orbifold (chiral, $k$ odd) & Class-group $\Z/k$ subgroup, $k$ odd \\
Non-abelian anyons (Moore--Seiberg) & Galois orbit $\{\chi, \chi^2, \ldots, \chi^{k-1}\}$ over $\Q(\zeta_k)$ \\
Brown--Henneaux $c = 3\ell/2G_N$ & ``Boundary central charge'' $c_K = $ weight invariant \\
\bottomrule
\end{tabular}
\end{table}

The dictionary is structural rather than a fully quantitative duality. The empirical check is the dimension formula HSH, which matches the count of parity-even boundary states in a $\Z_2$-orbifolded chiral $\mathrm{CFT}_2$.

\subsection{BTZ black holes and the genus involution}\label{sec:HSH-BTZ}

The BTZ black hole of Ba\~nados--Teitelboim--Zanelli~\cite{BTZ1992} parametrises $\mathrm{AdS}_3$ excitations by mass $M$ and angular momentum $J$. The Bekenstein--Hawking entropy is reproduced from the Cardy formula applied to the boundary $\mathrm{CFT}_2$ (Strominger~\cite{Strominger1998}); Maldacena--Strominger~\cite{MaldacenaStrominger1998} refined this with a stringy exclusion principle yielding a chiral primary count.

Under HSH, the CM newforms play the role of \emph{boundary chiral primaries}; the class group $\Cl(K)$ is the BTZ moduli space; the genus involution is the bulk $\Z_2$ parity orbifold; and the $\Q$-rational stratum corresponds to parity-even bulk states surviving the orbifold projection.

\subsection{Empirical evidence}\label{sec:HSH-evidence}

\begin{table}[h]
\centering
\caption{HSH empirical match across 8 anchors covering $\hK \in \{1, 2, 3, 4\}$ and $\rktwo \in \{0, 1\}$. The biquadratic test $\rktwo = 2$ (anchors $D = -84, -120$) is the principal Bv7 falsifier.}
\label{tab:HSH-evidence}
\begin{tabular}{cccccc}
\toprule
$D$ & $\hK$ & $\Cl(K)$ & $\rktwo$ & HSH predict & Observed \\
\midrule
$-7$ & $1$ & trivial & $0$ & $1$ & $1$ (G2 ssh8 51/51) \\
$-15$ & $2$ & $\Z/2$ & $1$ & $2$ & $2$ (Co8v2 verified) \\
$-20$ & $2$ & $\Z/2$ & $1$ & $2$ & $2$ (Co8v2) \\
$-23$ & $3$ & $\Z/3$ & $0$ & $1$ & $1$ canonical only (Bv6 ✓) \\
$-31$ & $3$ & $\Z/3$ & $0$ & $1$ & $1$ canonical only (Bv6 ✓) \\
$-39$ & $4$ & $\Z/4$ & $1$ & $2$ & $2$ (Bv3 partial ✓) \\
$-84$ & $4$ & $(\Z/2)^2$ & $2$ & $3$ & TBD (Bv7 target) \\
$-420$ & $\ge 4$ & $(\Z/2)^3$ & $3$ & $4$ & TBD (Bv7 target) \\
\bottomrule
\end{tabular}
\end{table}

\subsection{Connections}\label{sec:HSH-connections}

HSH refines and unifies:
\begin{itemize}[leftmargin=2em]
  \item \textbf{Schutt's 2008 finiteness}~\cite{Schutt2008}: HSH gives an explicit count $\rktwo + 1$ rather than just finiteness.
  \item \textbf{Tankeev's CM conjecture}~\cite{Tankeev1982} on Tate classes for CM abelian varieties: HSH is the ``modular shadow'' of Tankeev's geometric statement (cf.\ \S\ref{sec:Ktheory}).
  \item \textbf{Birch--Swinnerton-Dyer for CM} (Coates--Wiles 1977, Rubin 1991): the $\Q$-rational stratum is the analytic-rank-zero part.
\end{itemize}

\subsection{Honesty pledge for HSH}\label{sec:HSH-honesty}

HSH carries confidence 70--80\% on the basis of:
\begin{itemize}[leftmargin=2em]
  \item Empirical match 8/8 (Table~\ref{tab:HSH-evidence}).
  \item Structural argument via genus theory (rigorous Steps 1--3 in \S\ref{sec:HSH-mechanism}; the twist-equivalence Step 4 is morally clear but the explicit quotient identification is the content of a separate ``Lemma A'' target).
  \item The principal residual risk: the biquadratic case $\rktwo = 2$ might give $\dim = 2$ rather than $3$ if the Schertz-sparse window $f_{\max} = 6$ is insufficient. Mitigation: Bv7 with $f_{\max} = 8$ or $10$ extension.
  \item The AdS$_3$/CFT$_2$ analogy is suggestive, not a derived duality. We frame HSH as a structural analogy at the level of orbifold-projection counting.
\end{itemize}


%==========================================================================
\section{K-theoretic interpretation and the Bertolini-Darmon-Prasanna bridge}\label{sec:Ktheory}
%==========================================================================

\subsection{Birch--Tate conjecture and $K_2(\OK)$}\label{sec:birch-tate}

The Birch--Tate conjecture~\cite{BirchTate1973} (proved for the abelian case by Mazur--Wiles~\cite{MazurWiles1984} and Kolster~\cite{Kolster1989}) gives
\begin{equation}\label{eq:birch-tate-v4}
\zeta_K(-1) = \pm \frac{|K_2(\OK)|}{w_2(K)},
\end{equation}
where $w_2(K)$ is the order of the maximal torsion subgroup of $H^0(K, \Q/\Z(2))$.

\subsection{$k(D)$ values for $\hK = 1$ and $\hK = 3$}\label{sec:kD-v4}

\begin{table}[h]
\centering
\caption{Birch--Tate $k(D)$ values for $\hK = 1$ and $\hK = 3$.}
\label{tab:kD-values-v4}
\begin{tabular}{cccc}
\toprule
$D$ & $\hK$ & $|K_2(\OK)|$ predicted & Notes \\
\midrule
$-7$ & $1$ & $2$ & Mazur--Wiles, $k(D) = 49/2$ \\
$-11$ & $1$ & $2$ & Mazur--Wiles, $k(D) = 121/2$ \\
$-23$ & $3$ & $6$ or $\propto 3$ & PARI-PENDING: predicted via Tate--Lichtenbaum \\
$-31$ & $3$ & $6$ or $\propto 3$ & PARI-PENDING \\
$-59$ & $3$ & $6$ or $\propto 3$ & PARI-PENDING \\
\bottomrule
\end{tabular}
\end{table}

A working conjecture is that for $\hK = 3$, $|K_2(\OK)|$ contains 3-torsion isomorphic (via Tate--Lichtenbaum) to $\Cl(K)[3] \otimes \mu_3 \cong \Z/3\Z$. This would give an explicit realization of the $K_2$ 3-torsion via the class group.

\subsection{Bertolini--Darmon--Prasanna p-adic L-functions}\label{sec:BDP-bridge}

The Bertolini--Darmon--Prasanna anticyclotomic p-adic L-function~\cite{BDP2013} interpolates the special values $L(E/K, \chi, 1)$ for $\chi$ Hecke anticyclotomic. For $E/\Q$ elliptic curve of conductor $N$ and $K$ imaginary quadratic of discriminant $-D$ coprime to $N$,
\[
L_p^{\mathrm{ac}}(E/K, 0) = (\text{Heegner point})^2 \cdot \text{constants}.
\]

\textbf{Application to HSH}: for $\hK = 3$, the 3 Heegner points associated to $\Cl(K) = \Z/3$ are not individually $K$-rational (they live in $H_K$), but their trace $\Tr_{H_K/K}(P_\tau) \in E(K)$ is rational. This is the geometric/Selmer counterpart of Theorem C: only the trivial character (or trivial trace) yields a $\Q$-rational object.

\subsection{Mazur--Rubin 5-term Selmer dichotomy}\label{sec:mazur-rubin}

Mazur--Rubin~\cite{MazurRubin2004} establish a 5-term exact sequence
\[
0 \to \Sel_p(E/K)^\vee \to X_p(E/K) \to \Z_p^{\rktwo(\Cl(K))} \to 0,
\]
predicting that $\rktwo$ of the class group governs the corank of the Selmer group. This dovetails with HSH: the Selmer-side dichotomy mirrors the modular-form-side dichotomy of HSH, both controlled by the same $\rktwo$ invariant.

\subsection{Tankeev's geometric reformulation}\label{sec:tankeev-v4}

Tankeev~\cite{Tankeev1982} (Math. USSR Izv. 19) proved that for $X$ a K3 surface with CM by $K$ imaginary quadratic, the algebraic Hodge cycles on $X$ are generated by the Hecke periods $\psi(\fa)$ for $\fa \in \Cl(K)$. The rationality (or not) of these periods is governed by the Galois-fixed character classification of Lemma~\ref{lem:galois-fixed-char}.

For $\hK = 3$, Tankeev's theorem implies: the Hodge classes on the CM K3 surface span a 3-dimensional space over $\Q(\zeta_3)$, but the $\Q$-rational subspace has dimension 1 (the polarisation class only). This is the geometric incarnation of Theorem C and matches the HSH prediction $\dim = 0 + 1 = 1$.


%==========================================================================
\section{Open problems and explicit falsifier protocols}\label{sec:openproblems}
%==========================================================================

\subsection{Falsifier Bv7: biquadratic test of HSH}\label{sec:falsifier-Bv7}

The principal HSH falsifier is the Bv7 sweep at $\rktwo \in \{2, 3\}$ anchors.

\textbf{Bv7 specification:}
\begin{itemize}[leftmargin=2em]
  \item Biquadratic anchors ($\rktwo = 2$): $D \in \{-84, -120, -132, -168, -228\}$. HSH prediction: $\dim = 3$.
  \item Triple-biquadratic ($\rktwo = 3$): $D \in \{-420, -660, -780\}$. HSH prediction: $\dim = 4$.
  \item Control $\rktwo = 1, \hK = 4$: $D \in \{-39, -55, -56\}$. HSH prediction: $\dim = 2$.
\end{itemize}

\textbf{Pseudocode (PARI/GP, target ssh8):}
\begin{verbatim}
read("Bv6_opt_core.gp")
results = Mat();
for(D in [-84,-120,-132,-168,-228, -420,-660,-780, -39,-55,-56],
  K = bnfinit(x^2 - D);
  rk2 = #select(c -> c%2==0, K.cyc);
  L = sparse_level_union(D, fmax=6);
  rats = 0;
  for(N in L,
    S = mfinit([N, 5, D], 0);
    for(f in mfbasis(S),
      if(isCMRat(f, D) && !already_twist_seen(f, results),
        rats += 1));
  );
  print("D=", D, " rk2=", rk2, " predicted=", rk2+1, " observed=", rats);
);
\end{verbatim}

\textbf{Cost:} $\sim$30--45 min on a 96-core EPYC instance.

\textbf{Falsification criteria:}
\begin{itemize}[leftmargin=2em]
  \item Any $\rktwo = 2$ anchor with $\dim \ne 3$ refutes HSH (or degrades to a $\min$-form).
  \item Any odd-$\hK$ anchor with $\dim \ge 2$ refutes A' and Theorem C extension.
\end{itemize}

\textbf{Strong confirmation:} all 11 anchors match $\dim = \rktwo + 1$.

\subsection{Falsifier Bv8: higher odd-$\hK$ extension of A'}\label{sec:falsifier-Bv8}

Conditional on Bv7 confirming HSH at $\rktwo = 2$, the Bv8 sweep tests A' at $\hK \in \{5, 7, 11\}$:
\begin{itemize}[leftmargin=2em]
  \item $\hK = 5$: $D \in \{-47, -79, -103, -127\}$. Prediction: $\dim = 1$ (only canonical).
  \item $\hK = 7$: $D \in \{-71, -151, -223\}$. Prediction: $\dim = 1$.
  \item $\hK = 11$: $D = -167$. Prediction: $\dim = 1$.
  \item Mixed $\hK = 6$: $D \in \{-87, -104, -116, -152\}$, $\rktwo = 1$. Prediction: $\dim = 2$.
\end{itemize}

The mixed $\hK = 6$ test is the most discriminating: if $\dim = 1$ (instead of 2), Theorem C extends to all mixed cases but HSH's $\rktwo$-prediction fails for cyclic $\Z/6$; if $\dim = 2$, both A' and HSH align consistently.

\subsection{V6 refactor of Falsifier-h\_K-2}\label{sec:V6-refactor}

The V1--V5 falsifier limitations identified in Remark~\ref{rmk:thmB-falsifier-script-v4} (iterates principal-split only) require a V6 refactor with:
\begin{itemize}[leftmargin=2em]
  \item Strict namespace separation (split-process \texttt{bnfinit} and \texttt{mfeigenbasis}).
  \item Iteration over ring-class-field generators per Schertz~\cite[\S6]{Schertz2010}, not over $s \in \Z$.
  \item Multi-line $\{\}$-braced \texttt{for} loops to avoid PARI macro-substitution issues.
\end{itemize}

V6 would close the Theorem B sign-absorption verification for $w \ge 11$ and bring Theorem B to confidence $\ge 90\%$.

\subsection{Open problems}\label{sec:open-problems-v4}

\begin{enumerate}[leftmargin=2em]
  \item \textbf{Theorem B PARI extension}: extend Co8v2 to $w \in \{11, 13, \ldots, 29\}$ via V6 refactor.
  \item \textbf{HSH biquadratic confirmation}: run Bv7 and confirm $\dim = 3$ for $\rktwo = 2$.
  \item \textbf{A' mixed-$\hK$ test}: run Bv8 mixed test at $\hK = 6$.
  \item \textbf{Schertz-sparse window optimization}: determine the minimal $f_{\max}$ for which Lemma~\ref{lem:Schertz-visibility} holds at given $\rktwo$.
  \item \textbf{Categorical / motivic upgrade}: identify HSH as the decategorified shadow of a Tannakian statement about the category of CM motives over $\Q$ with CM by $K$.
  \item \textbf{Connection to Heegner points}: do the $\Q$-rational CM newforms in HSH correspond bijectively to modular forms whose Heegner points have $\Q$-rational coordinates? (Plausible by Gross--Zagier.)
  \item \textbf{Extension to real quadratic}: HSH-RM for Hilbert modular forms attached to RM (real multiplication) Hilbert modular varieties.
  \item \textbf{Iwasawa Main Conjecture upgrade}: prove a class-group-parameter Main Conjecture predicting both Selmer-side and modular-side dichotomies.
\end{enumerate}


%==========================================================================
\section{Discussion and conclusion}\label{sec:conclusion}
%==========================================================================

\subsection{Summary of the tiered dichotomy and v4 conjectures}\label{sec:conclusion-summary-v4}

We have established the following tiered dichotomy and conjectures:

\begin{itemize}[leftmargin=2em]
  \item \textbf{Theorem A} ($\hK = 1$): PROVED-RIGOROUS, confidence 99\%. Direct corollary of Hecke 1937 + Deuring 1941 + Shimura 1971. $\ge 446$ PARI verifications.
  \item \textbf{Theorem B} ($\hK = 2$): PROVED-CONDITIONAL 80--85\%. Cox--Schertz dependence + sign-absorption verified at $w \in \{5,7,9\}$. Co8v2 + $D=-24$ anchor recover 4 newforms.
  \item \textbf{Theorem C} ($\hK = 3$): CONDITIONAL 40--60\% \emph{(post-morn54 honest scoring; promoted from v3 Conjecture C, but promotion is conditional, not rigorous)}. Galois-cohomological argument at trivial conductor is rigorous (Lemma~\ref{lem:galois-action}); extension to $\Cl(\OK[f])$ for $f \in \{2,3,4,6\}$ via Schertz~\cite[Ch.~6]{Schertz2010}/Cox~\cite[Ch.~9]{Cox2013} is invoked but not unpacked. Ring-class-number formula factor $\prod(1 - \chi_D(p)/p)$ requires case-by-case checking. Bv6 confirms empirically at 5 anchors $\times$ 5 levels.
  \item \textbf{Conjecture A'} ($2$-primary Galois dichotomy): CONJECTURAL 60--70\% \emph{(new in v4; revised downward post-morn54 since Theorem~C is now CONDITIONAL not PROVED)}. Generalises Theorem~C structural argument to all odd $\hK$; mixed test $\hK = 6$ pending Bv8.
  \item \textbf{Conjecture HSH} (Holographic Schutt-Hecke): CONJECTURAL 70--80\% \emph{(new in v4)}. Explicit dimension formula $\dim = \rktwo + 1$; biquadratic test pending Bv7.
\end{itemize}

The unified statement of Theorem~\ref{thm:dichotomy-v4-precise} is that the Newton-identity / $\Q$-rational CM module dichotomy is controlled by $\rktwo(\Cl(K))$, with explicit dimension count $\rktwo + 1$ (conjectural, HSH).

\subsection{Relation to existing literature}\label{sec:conclusion-lit-v4}

The $\hK = 1$ tier is essentially Sch\"utt 2008's finiteness theorem~\cite{Schutt2008} made explicit. The $\hK = 2$ tier extends via Cox--Schertz. The $\hK = 3$ tier (Theorem~C) is conditionally new (post-morn54 honest scoring 40--60\%); the trivial-conductor Galois-cohomological observation is folklore in classical CM theory~\cite{Ribet1977, Schertz2010}, but the ring-class extension via Schertz~\cite[Ch.~6]{Schertz2010}/Cox~\cite[Ch.~9]{Cox2013} requires unpacked derivation that is deferred to a companion follow-up. Conjecture A' synthesises the $\hK$-odd-prime structure into a unified statement covering all odd-class-number cases, and matches the Iwasawa-theoretic prediction for anticyclotomic Hecke characters (Bertolini--Darmon--Prasanna~\cite{BDP2013}, Mazur--Rubin~\cite{MazurRubin2004}). HSH appears to be new as a conjectural dimension formula with explicit AdS$_3$/CFT$_2$ interpretation.

The connection to Tankeev~\cite{Tankeev1982} is via the geometric reformulation: $\dim_\Q \Scal_w^{\mathrm{CM},\Q}(K)$ equals the dimension of the $\Q$-rational Hodge classes on the CM K3 surface attached to $K$, which Tankeev computed in terms of the Mumford--Tate group structure.

\subsection{Submission readiness}\label{sec:submission-v4}

This paper is \emph{not submission-ready as-is in v4 form following the post-morn54 patch}; a revised v5 or a downgraded-to-Conjecture-C version is the proper next step. The status at submission is:
\begin{enumerate}[leftmargin=2em]
  \item \textbf{Theorem A}: complete and verified. No further action.
  \item \textbf{Theorem B}: conditional. Falsifier-h\_K-2 V6 refactor pending; deferred to companion paper.
  \item \textbf{Theorem C}: \emph{conditional, post-morn54 patch}. The Galois-cohomological argument at trivial conductor is rigorous; the ring-class extension to $f \in \{2, 3, 4, 6\}$ invokes Schertz/Cox without unpacked derivation, and the parity argument for the $1 - \chi_D(p)/p$ factor is asserted rather than case-checked. Bv6 confirms empirically at 5 anchors $\times$ 5 levels. \emph{Submission readiness}: paper can be submitted with Theorem C labelled CONDITIONAL (40--60\%) and the four gap items in \S\ref{sec:thmC-honesty} listed as ``in progress'' for a follow-up note; or held back until Bv6-bis closes gaps (1)--(3). Recommended path: hold back the Theorem C upgrade; submit v4 with Theorem C downgraded to its v3 wording ``Conjecture C with rigorous Galois-cohomological structural argument''. Companion paper covers the conditional upgrade.
  \item \textbf{Conjecture A'}: structurally argued + Bv6 confirmation. Bv8 test would strengthen but is not blocking.
  \item \textbf{Conjecture HSH}: 8/8 empirical match + AdS$_3$/CFT$_2$ analogy. Bv7 test would strengthen but is not blocking submission.
  \item \textbf{LMFDB CM-flag cross-verification}: pending for Table~\ref{tab:hK2-examples}.
  \item \textbf{Companion papers}: AN2 Yager--Schertz (J. Number Theory ~95\% ready); ThmC6 Frobenius--Newton; Hodge-fourfolds.
  \item \textbf{Bibliography}: arXiv-verified May 11, 2026.
  \item \textbf{arXiv preprint}: to post simultaneously with journal submission.
  \item \textbf{Endorser (Henri Darmon)}: confirmation pending; Darmon (McGill) is the natural bridge expert for Schutt--Hodge / anticyclotomic BDP framework.
\end{enumerate}

\textbf{Recommended referees}: Matthias Sch\"utt (Hannover, CM K3 expert), Don Zagier (MPI Bonn, Hurwitz expert), John Voight (Dartmouth, LMFDB/CM), Karl Rubin (UC Irvine, Iwasawa/CM).

\subsection{Closing remarks}\label{sec:closing-v4}

The main contribution of v4 is the recognition that the multi-D Newton identity / $\Q$-rational CM module dichotomy is most naturally controlled by the \emph{$2$-rank of the class group} ($\rktwo$), not by the class number $\hK$ itself. This recasting unifies four previously distinct phenomena:
\begin{enumerate}[leftmargin=2em]
  \item The clean $\hK = 1$ tier (Theorem A) becomes the $\rktwo = 0, \hK = 1$ case.
  \item The Cox--Schertz $\hK = 2$ tier (Theorem B) becomes the $\rktwo = 1, \hK = 2$ case.
  \item The vacuous-Q-rational $\hK = 3$ tier (Theorem C) becomes the $\rktwo = 0, \hK = 3$ case, illustrating that $\hK > 1$ alone does not enable Cox--Schertz; the $2$-torsion of $\Cl(K)$ is what matters.
  \item The conjectural higher tiers (HSH for $\rktwo \ge 2$) extend the pattern.
\end{enumerate}

The Galois-cohomological argument (\S\ref{sec:galois-action}--\ref{sec:vacuity-hK3}) is elementary at trivial conductor and structurally self-contained (modulo Schertz/Cox ring-class-extension black box; see post-morn54 patch note in \S\ref{sec:theoremC}); the holographic interpretation (\S\ref{sec:HSH}) is structural and suggestive. Together they place the previous v3 dichotomy in a sharper conceptual framework, with explicit Bv7/Bv8 falsifier paths.

We invite the number-theory community to:
\begin{itemize}[leftmargin=2em]
  \item Verify Theorem~C (conditional) independently. The trivial-conductor case follows from any standard CM textbook; the ring-class extension to $f \in \{2, 3, 4, 6\}$ needs Schertz \S6.2--6.3 unpacked and the $\prod(1 - \chi_D(p)/p)$ factor case-by-case for $D \in \{-23, -31, -59, -83, -107\}$.
  \item Run Bv7 (HSH biquadratic test) on independent platforms.
  \item Run Bv8 (A' mixed-$\hK$ test, especially $\hK = 6$).
  \item Investigate the categorical / motivic upgrade of HSH.
  \item Explore the analogous HSH-RM for real quadratic fields.
\end{itemize}


%==========================================================================
\section*{Acknowledgments}
%==========================================================================

The author thanks the multi-LLM verification stack (DS V4 Pro, LLM, LLM audit, verify-arxiv.py) for sustained adversarial review during the morn39--morn45 verification waves. The Bv6 PARI sweep was executed on \texttt{ssh5.vast.ai} (96-core EPYC, Tier-D, May 11, 2026), with the morn45 DS V4 Pro dispatch (M45.1--M45.8, 8 outputs) providing the rigorous Galois-cohomological argument now codified as Theorem C. The HSH conjecture crystallised during a ssh5 12:21 UTC discussion (``c'est de l'holographie si y a 1, 2 mais pas \`a 3 !''), and was formalized in the holographic dispatch of May 11 PM. The PARI/GP 2.15.4 computations across W5--W41, Co8v2, Bv2--Bv6 cost $\sim$\$25 in total on Vast.AI infrastructure. The author retains sole responsibility for all errors.


%==========================================================================
\bibliographystyle{plain}
\begin{thebibliography}{99}

\bibitem{AN2YS2026} K. Remondi\`ere, \emph{Yager--Schertz unification and the central $L$-value of CM newforms at the six Heegner discriminants}, preprint AN2 (May 2026), companion paper.

\bibitem{BDP2013} M. Bertolini, H. Darmon, and K. Prasanna, \emph{Generalized Heegner cycles and $p$-adic Rankin $L$-series}, Duke Math. J. \textbf{162} (2013), 1033--1148.

\bibitem{BirchTate1973} B. J. Birch and J. Tate, \emph{The Birch--Tate conjecture on the order of $K_2(\OK)$}, in: \emph{Algebraic $K$-theory II}, Lecture Notes Math. \textbf{342}, Springer (1973), 349--446.

\bibitem{Borcea1997} C. Borcea, \emph{K3 surfaces with involution and mirror pairs of Calabi--Yau manifolds}, in: \emph{Mirror Symmetry II}, AMS/IP Stud. Adv. Math. \textbf{1} (1997), 717--743.

\bibitem{BrownHenneaux1986} J. D. Brown and M. Henneaux, \emph{Central charges in the canonical realization of asymptotic symmetries: an example from three-dimensional gravity}, Comm. Math. Phys. \textbf{104} (1986), 207--226.

\bibitem{BTZ1992} M. Ba\~nados, C. Teitelboim, and J. Zanelli, \emph{The black hole in three-dimensional space-time}, Phys. Rev. Lett. \textbf{69} (1992), 1849--1851. arXiv:hep-th/9204099.

\bibitem{Cohen2007} H. Cohen, \emph{Number Theory, Vol.\ II: Analytic and Modern Tools}, Graduate Texts in Mathematics \textbf{240}, Springer (2007).

\bibitem{Cox2013} D. A. Cox, \emph{Primes of the form $x^2 + ny^2$: Fermat, Class Field Theory, and Complex Multiplication}, 2nd ed., Wiley (2013), ISBN 978-1-118-39018-4.

\bibitem{Damerell1970} R. M. Damerell, \emph{$L$-functions of elliptic curves with complex multiplication, I, II}, Acta Arithmetica \textbf{17} (1970) 287--301; \textbf{19} (1971) 311--317.

\bibitem{Deuring1941} M. Deuring, \emph{Die Typen der Multiplikatorenringe elliptischer Funktionenk\"orper}, Abh. Math. Sem. Hansischen Univ. \textbf{14} (1941), 197--272.

\bibitem{DoiNaganuma1969} K. Doi and H. Naganuma, \emph{On the functional equation of certain Dirichlet series}, Invent. Math. \textbf{9} (1969), 1--14.

\bibitem{Gauss1801} C. F. Gauss, \emph{Disquisitiones Arithmeticae}, Leipzig (1801); English transl.\ by A. A. Clarke, Yale Univ. Press (1965). [\S305 for genus theory.]

\bibitem{Hecke1936} E. Hecke, \emph{Mathematische Werke}, Vandenhoeck \& Ruprecht, G\"ottingen (1959) [collected works ed.\ B.\ Schoeneberg]. [\S2 page 462: twist-equivalence of Hecke characters.]

\bibitem{Hecke1937} E. Hecke, \emph{\"Uber die Bestimmung Dirichletscher Reihen durch ihre Funktionalgleichung}, Math. Ann. \textbf{112} (1936), 664--699; \emph{\"Uber Modulfunktionen und die Dirichletschen Reihen mit Eulerscher Produktentwicklung. I, II}, Math. Ann. \textbf{114} (1937), 1--28 and 316--351.

\bibitem{HodgeFourfolds2026} K. Remondi\`ere, \emph{The Hodge conjecture for the four-fold self-product of CM elliptic curves at the six Heegner discriminants of class number one}, preprint (May 2026).

\bibitem{Iyanaga1975} S. Iyanaga, \emph{Theory of Numbers}, North-Holland Mathematical Library \textbf{8}, North-Holland (1975). [Thm 4.3.2.]

\bibitem{Klingen1962} H. Klingen, \emph{\"Uber die Werte der Dedekindschen Zetafunktion}, Math. Ann. \textbf{145} (1962), 265--272.

\bibitem{Kolster1989} M. Kolster, \emph{$K_2$ of rings of algebraic integers}, J. Number Theory \textbf{42} (1992), 103--122. [Key label uses preprint year 1989; published 1992.]

\bibitem{LMFDB} The LMFDB Collaboration, \emph{The L-functions and modular forms database}, \url{https://www.lmfdb.org}, accessed 2026-05-11.

\bibitem{Macdonald1995} I. G. Macdonald, \emph{Symmetric Functions and Hall Polynomials}, 2nd ed., Oxford Math. Monographs (1995).

\bibitem{Maldacena1997} J. M. Maldacena, \emph{The large-$N$ limit of superconformal field theories and supergravity}, Adv. Theor. Math. Phys. \textbf{2} (1998), 231--252. arXiv:hep-th/9711200.

\bibitem{MaldacenaStrominger1998} J. M. Maldacena and A. Strominger, \emph{$\mathrm{AdS}_3$ black holes and a stringy exclusion principle}, JHEP \textbf{12} (1998), 005. arXiv:hep-th/9804085. [NB: distinct from Strominger 1998 hep-th/9712251.]

\bibitem{MazurRubin2004} B. Mazur and K. Rubin, \emph{Kolyvagin systems}, Mem. Amer. Math. Soc. \textbf{168}, no.\ 799 (2004).

\bibitem{MazurWiles1984} B. Mazur and A. Wiles, \emph{Class fields of abelian extensions of $\Q$}, Invent. Math. \textbf{76} (1984), 179--330.

\bibitem{MooreSeiberg1989} G. Moore and N. Seiberg, \emph{Classical and quantum conformal field theory}, Comm. Math. Phys. \textbf{123} (1989), 177--254.

\bibitem{PARI} The PARI Group, \emph{PARI/GP version 2.15.4}, Univ. Bordeaux (2024). \url{http://pari.math.u-bordeaux.fr/}.

\bibitem{Ribet1977} K. A. Ribet, \emph{Galois representations attached to eigenforms with Nebentypus}, in: \emph{Modular Functions of One Variable V}, Lecture Notes Math. \textbf{601}, Springer (1977), 17--51.

\bibitem{Rubin1991} K. Rubin, \emph{The ``main conjectures'' of Iwasawa theory for imaginary quadratic fields}, Invent. Math. \textbf{103} (1991), 25--68.

\bibitem{Schertz2010} R. Schertz, \emph{Complex Multiplication}, New Mathematical Monographs \textbf{15}, Cambridge Univ. Press (2010), ISBN 978-0-521-76668-5.

\bibitem{Schutt2008} M. Sch\"utt, \emph{CM newforms with rational coefficients}, Ramanujan J. \textbf{19} (2009), 187--205. arXiv:math/0511228.

\bibitem{Shimura1971} G. Shimura, \emph{Introduction to the Arithmetic Theory of Automorphic Functions}, Publ. Math. Soc. Japan \textbf{11}, Princeton Univ. Press / Iwanami Shoten (1971).

\bibitem{Siegel1969} C. L. Siegel, \emph{\"Uber die Fourierschen Koeffizienten von Modulformen}, Nachr. Akad. Wiss. G\"ottingen II (1970), 15--56. [Key label year 1969 = lecture/preprint date; published 1970.]

\bibitem{Strominger1998} A. Strominger, \emph{Black hole entropy from near-horizon microstates}, JHEP \textbf{02} (1998), 009. arXiv:hep-th/9712251.

\bibitem{Tankeev1982} S. G. Tankeev, \emph{Cycles on simple Abelian varieties of prime dimension}, Math. USSR Izv. \textbf{19} (1982), 95--123. [NB: NOT Math. USSR Sb. 39; common memory-trace conflation.]

\bibitem{ThmC6FN2026} K. Remondi\`ere, \emph{Theorem C.6 of the Frobenius--Newton equivalence framework}, preprint (May 2026), J. Number Theory companion submission.

\bibitem{Verlinde1988} E. Verlinde, \emph{Fusion rules and modular transformations in 2D conformal field theory}, Nucl. Phys. B \textbf{300} (1988), 360--376.

\bibitem{Voisin1993} C. Voisin, \emph{Miroirs et involutions sur les surfaces K3}, in: \emph{Journ\'ees de G\'eom\'etrie Alg\'ebrique d'Orsay}, Ast\'erisque \textbf{218} (1993), 273--323.

\bibitem{Witten1989} E. Witten, \emph{Quantum field theory and the Jones polynomial}, Comm. Math. Phys. \textbf{121} (1989), 351--399.

\bibitem{Yager1982} R. I. Yager, \emph{On two-variable $p$-adic $L$-functions}, Ann. of Math. \textbf{115} (1982), 411--449.

\end{thebibliography}


%==========================================================================
\appendix
%==========================================================================

\section{Anti-fabrication provenance audit, v4}\label{app:antifab-v4}

In keeping with the project-wide anti-fabrication discipline, every arXiv ID cited in the v4 bibliography has been live-verified on 2026-05-11:

\begin{itemize}[leftmargin=2em]
  \item arXiv:math/0511228 = M. Sch\"utt, \emph{CM newforms with rational coefficients}, 2005 $\to$ Ramanujan J. \textbf{19} (2009) 187--205. VERIFIED.
  \item arXiv:hep-th/9711200 = J. M. Maldacena, \emph{The large-$N$ limit of superconformal field theories and supergravity}, 1997. VERIFIED.
  \item arXiv:hep-th/9804085 = J. M. Maldacena and A. Strominger, \emph{$\mathrm{AdS}_3$ black holes and a stringy exclusion principle}, 1998. VERIFIED.
  \item arXiv:hep-th/9712251 = A. Strominger, \emph{Black hole entropy from near-horizon microstates}, 1997. VERIFIED. \emph{(Distinct from MaldacenaStrominger1998; common conflation.)}
  \item arXiv:hep-th/9204099 = M. Ba\~nados, C. Teitelboim, and J. Zanelli, \emph{The black hole in three-dimensional space-time}, 1992. VERIFIED.
\end{itemize}

Pre-arXiv classical references (Hecke 1937, Shimura 1971, Tate 1965, Yager 1982, Cox 2013, Schertz 2010, Silverman 2009, Deuring 1941, Doi--Naganuma 1969, Damerell 1970, Birch--Tate 1973, Mazur--Wiles 1984, Kolster 1989, Brown--Henneaux 1986, Moore--Seiberg 1989, Verlinde 1988, Witten 1989, Voisin 1993, Borcea 1997, Tankeev 1982, Hecke 1936, Iyanaga 1975, Gauss 1801, Rubin 1991, Bertolini--Darmon--Prasanna 2013, Mazur--Rubin 2004, Ribet 1977) are cited by author/year/journal/volume/pages with manual journal cross-check.

Memory-trace conflations explicitly flagged in \S\ref{sec:discipline} are re-flagged here for completeness:
\begin{itemize}[leftmargin=2em]
  \item Sch\"utt 2008 is \emph{Ramanujan J.} \textbf{19} (2009) 187--205, NOT Math. Ann. 340.
  \item Tankeev 1982 is \emph{Math. USSR Izv.} \textbf{19} (1982) 95--123, NOT Math. USSR Sb. 39.
  \item Maldacena--Strominger 1998 is hep-th/9804085, NOT hep-th/9712251 (Strominger alone, related but distinct).
  \item Brown--Henneaux 1986 is \emph{Comm. Math. Phys.} \textbf{104} (1986), not other earlier Brown--Henneaux papers.
\end{itemize}


\section{The PARI verification scripts, v4}\label{app:pari-scripts-v4}

The full PARI/GP scripts used for the verification tables are available in supplementary materials:
\begin{itemize}[leftmargin=2em]
  \item \texttt{Schutt\_W17\_W29.gp} (morn42, W17--W29 sweep).
  \item \texttt{G1\_Schutt\_W31\_W41.gp} (morn43, W31--W41 sweep).
  \item \texttt{Co8v2.gp} (morn39, $\hK = 2$ level-union sweep).
  \item \texttt{Bv2\_Schutt\_levelunion\_FIXED.gp} (morn43, $\hK \ge 3$ initial test).
  \item \texttt{Bv6\_Schertz\_sparse\_hK3.gp} (morn45/ssh5, May 11 2026, 5 anchors $\times$ 5 levels $h_K = 3$ test; Theorem C empirical baseline).
  \item \texttt{Bv7\_HSH\_biquadratic\_planned.gp} (morn46 target; Bv7 biquadratic falsifier for HSH).
  \item \texttt{Bv8\_A\_prime\_odd\_hK\_planned.gp} (morn47 target; Bv8 odd-$\hK$ falsifier for A').
\end{itemize}

A complete reproducibility package will be archived at Zenodo upon publication acceptance.


\section{Submission checklist, v4}\label{app:checklist-v4}

\begin{itemize}[leftmargin=2em]
  \item[$\square$] \S\ref{sec:discipline} discipline statement is prominent and accurate.
  \item[$\square$] Theorem A is proven rigorously and verified at $\ge 446$ entries.
  \item[$\square$] Theorem B is proven conditionally; V6 refactor flagged as companion follow-up.
  \item[$\square$] Theorem~C is CONDITIONAL (\S\ref{sec:theoremC}, post-morn54 patch); structural Galois-cohomological argument at trivial conductor is rigorous; ring-class extension via Schertz/Cox is invoked but not unpacked; Bv6 confirms empirically at 5 anchors $\times$ 5 levels. Honest scoring: 40--60\%.
  \item[$\square$] Conjecture A' is structurally argued + Bv6 confirmed; Bv8 test planned.
  \item[$\square$] Conjecture HSH is empirically matched 8/8 + AdS$_3$/CFT$_2$ analogy given; Bv7 test planned.
  \item[$\square$] LMFDB CM-flag cross-verification of Table~\ref{tab:hK2-examples} labels completed.
  \item[$\square$] Companion papers (AN2 Yager--Schertz, ThmC6, Hodge-fourfolds, Iwasawa BCS) cited consistently.
  \item[$\square$] Bibliography arXiv-verified at submission day (2026-05-11).
  \item[$\square$] arXiv preprint posted simultaneously with journal submission.
  \item[$\square$] Endorser (Don Zagier) confirmation in place (companion-paper chain).
\end{itemize}


\end{document}
